\documentclass[8pt]{beamer}
\usetheme{Madrid}
\usecolortheme{default}
\usepackage{amsmath, amssymb, graphicx, array, multirow, booktabs, xcolor, tikz}
\usetikzlibrary{shapes,arrows,positioning,fit,backgrounds}

% Compact font settings
\setbeamerfont{title}{size=\Large}
\setbeamerfont{subtitle}{size=\small}
\setbeamerfont{author}{size=\scriptsize}
\setbeamerfont{date}{size=\scriptsize}
\setbeamerfont{frametitle}{size=\large}
\setbeamerfont{framesubtitle}{size=\normalsize}
\setbeamerfont{enumerate item}{size=\footnotesize}
\setbeamerfont{itemize item}{size=\footnotesize}
\setbeamerfont{block title}{size=\footnotesize}

% Compact spacing
\setlength{\itemsep}{0.08ex}
\setlength{\parskip}{0.08ex}
\setlength{\leftmargini}{0.3cm}
\setbeamersize{text margin left=3mm, text margin right=3mm}
\addtobeamertemplate{frame begin}{\vspace{-0.4cm}}{}

% Colors
\definecolor{myblue}{RGB}{0,0,255}
\definecolor{myred}{RGB}{255,0,0}
\definecolor{mygreen}{RGB}{0,128,0}
\definecolor{myorange}{RGB}{255,165,0}
\definecolor{mypurple}{RGB}{128,0,128}
\definecolor{mygray}{RGB}{128,128,128}
\definecolor{mygold}{RGB}{255,215,0}

\title{Data and AI Model Governance}
\subtitle{100 Questions: Data Governance, Model Risk Management, Validation, GenAI Governance}
\author{GARP RAI Exam Preparation}
\date{}

\begin{document}
	
	\begin{frame}
		\titlepage
	\end{frame}
	
	% ============================================================
	% SECTION 1: INTRODUCTION TO DATA AND MODEL GOVERNANCE - QUESTIONS 1-5
	% ============================================================
	
	\begin{frame}{Q1: Model Governance Definition - Tutorial}
		\fontsize{6.5}{7.5}\selectfont
		\textbf{QUESTION: What is the primary purpose of model governance in financial institutions?}
		
		\vspace{0.1cm}
		\textbf{OPTIONS:}
		\begin{enumerate}
			\item To increase the number of models used
			\item To ensure effective and prudent risk management across the entire model lifecycle
			\item To eliminate all model risks completely
			\item To reduce the cost of model development
		\end{enumerate}
		
		\vspace{0.1cm}
		\textcolor{myblue}{\textbf{ANSWER: B}}
		
		\vspace{0.1cm}
		\textbf{DETAILED EXPLANATION:}
		\begin{itemize}
			\item \textbf{Why B is correct:} Sound governance of models is essential for effective and prudent risk management. Model governance should exist across the entire model lifecycle - from development through performance monitoring and decommissioning. Without it, a range of outcomes and decisions may occur that can prove detrimental.
			\item \textbf{Why A is wrong:} Increasing the number of models is not the purpose of governance
			\item \textbf{Why C is wrong:} Complete elimination of all model risks is impossible - models are simplifications of reality
			\item \textbf{Why D is wrong:} While good governance may reduce costs, that is not its primary purpose
		\end{itemize}
		
		\textcolor{myred}{\textbf{EXAM TRAP:}} Model governance is about \textcolor{myblue}{risk management across the entire lifecycle} - not just development!
	\end{frame}
	
	\begin{frame}{Q2: Data Governance Definition - Tutorial}
		\fontsize{6.5}{7.5}\selectfont
		\textbf{QUESTION: According to the Data Governance Institute, what is data governance?}
		
		\vspace{0.1cm}
		\textbf{OPTIONS:}
		\begin{enumerate}
			\item A method for collecting as much data as possible
			\item A system of decision rights and accountabilities for information-related processes
			\item A software tool for data storage
			\item A regulatory requirement that only applies to banks
		\end{enumerate}
		
		\vspace{0.1cm}
		\textcolor{myblue}{\textbf{ANSWER: B}}
		
		\vspace{0.1cm}
		\textbf{DETAILED EXPLANATION:}
		\begin{itemize}
			\item \textbf{Why B is correct:} The Data Governance Institute defines data governance as "a system of decision rights and accountabilities for information-related processes, executed according to agreed-upon models which describe who can take what actions with what information, and when, under what circumstances, using what methods."
			\item \textbf{Why A is wrong:} Data governance is about quality and control, not quantity
			\item \textbf{Why C is wrong:} Governance encompasses people, processes, and technologies - not just software
			\item \textbf{Why D is wrong:} Data governance applies across industries, not just banking
		\end{itemize}
		
		\textcolor{myred}{\textbf{EXAM TRAP:}} Data governance is about \textcolor{myblue}{decision rights and accountabilities} - not just collecting data!
	\end{frame}
	
	\begin{frame}{Q3: AI/ML Governance Importance - Tutorial}
		\fontsize{6.5}{7.5}\selectfont
		\textbf{QUESTION: Why is sound governance of AI/ML models potentially even more important than for traditional models?}
		
		\vspace{0.1cm}
		\textbf{OPTIONS:}
		\begin{enumerate}
			\item AI/ML models are always simpler than traditional models
			\item The opacity of AI/ML models makes proper governance of data and outcomes particularly relevant
			\item AI/ML models never need validation
			\item Traditional models have no risks
		\end{enumerate}
		
		\vspace{0.1cm}
		\textcolor{myblue}{\textbf{ANSWER: B}}
		
		\vspace{0.1cm}
		\textbf{DETAILED EXPLANATION:}
		\begin{itemize}
			\item \textbf{Why B is correct:} Sound governance of AI/ML models is even more important as it helps ensure technological advances do not overshadow the need for fundamental challenges. The opacity of AI/ML models makes the proper governance of the data used to train these models particularly relevant.
			\item \textbf{Why A is wrong:} AI/ML models are often more complex than traditional models
			\item \textbf{Why C is wrong:} AI/ML models require rigorous validation
			\item \textbf{Why D is wrong:} Traditional models have significant risks that need governance
		\end{itemize}
		
		\textcolor{myred}{\textbf{EXAM TRAP:}} AI/ML governance is \textcolor{myblue}{especially important due to opacity and data dependencies}!
	\end{frame}
	
	\begin{frame}{Q4: Model Lifecycle - Tutorial}
		\fontsize{6.5}{7.5}\selectfont
		\textbf{QUESTION: What phases should model governance cover?}
		
		\vspace{0.1cm}
		\textbf{OPTIONS:}
		\begin{enumerate}
			\item Only the development phase
			\item Only the deployment phase
			\item The entire model lifecycle from development through performance monitoring and decommissioning
			\item Only the validation phase
		\end{enumerate}
		
		\vspace{0.1cm}
		\textcolor{myblue}{\textbf{ANSWER: C}}
		
		\vspace{0.1cm}
		\textbf{DETAILED EXPLANATION:}
		\begin{itemize}
			\item \textbf{Why C is correct:} Model governance should exist across the entire model lifecycle - from model development through performance monitoring and ultimately decommissioning. This ensures comprehensive oversight at all stages.
			\item \textbf{Why A is wrong:} Governance must extend beyond development
			\item \textbf{Why B is wrong:} Governance must include pre-deployment phases
			\item \textbf{Why D is wrong:} Validation is part of governance, not the whole picture
		\end{itemize}
		
		\textcolor{myred}{\textbf{EXAM TRAP:}} Model governance covers \textcolor{myblue}{the entire lifecycle} - not just one phase!
	\end{frame}
	
	\begin{frame}{Q5: Regulatory Framework - Tutorial}
		\fontsize{6.5}{7.5}\selectfont
		\textbf{QUESTION: Which regulatory document provides key guidance on model risk management in the US financial sector?}
		
		\vspace{0.1cm}
		\textbf{OPTIONS:}
		\begin{enumerate}
			\item GDPR (General Data Protection Regulation)
			\item SR 11-7 (Supervisory Letter on Model Risk Management)
			\item Basel III
			\item Sarbanes-Oxley Act
		\end{enumerate}
		
		\vspace{0.1cm}
		\textcolor{myblue}{\textbf{ANSWER: B}}
		
		\vspace{0.1cm}
		\textbf{DETAILED EXPLANATION:}
		\begin{itemize}
			\item \textbf{Why B is correct:} The Federal Reserve's Supervisory Letter SR 11-7 from 2011 established regulatory standards concerning model governance. Although general in nature, these standards are relevant to AI/ML models.
			\item \textbf{Why A is wrong:} GDPR is European privacy regulation, not US model risk guidance
			\item \textbf{Why C is wrong:} Basel III is international banking capital regulation
			\item \textbf{Why D is wrong:} Sarbanes-Oxley is corporate governance and financial reporting legislation
		\end{itemize}
		
		\textcolor{myred}{\textbf{EXAM TRAP:}} SR 11-7 is the \textcolor{myblue}{key US regulatory guidance} for model risk management!
	\end{frame}
	
	% ============================================================
	% SECTION 2: DATA GOVERNANCE - QUESTIONS 6-20
	% ============================================================
	
	\begin{frame}{Q6: Data Strategy - Tutorial}
		\fontsize{6.5}{7.5}\selectfont
		\textbf{QUESTION: What should a comprehensive data strategy include?}
		
		\vspace{0.1cm}
		\textbf{OPTIONS:}
		\begin{enumerate}
			\item Only data collection methods
			\item Clear data-governance policy, strategic vision, operational practices, and authorization levels
			\item Only data storage procedures
			\item Only data deletion policies
		\end{enumerate}
		
		\vspace{0.1cm}
		\textcolor{myblue}{\textbf{ANSWER: B}}
		
		\vspace{0.1cm}
		\textbf{DETAILED EXPLANATION:}
		\begin{itemize}
			\item \textbf{Why B is correct:} Data procurement and use requires a clear and approved data-governance policy that is strategic and outlines operational practices. This should include a level of authorization for non-publicly available personal information, specification of whether data is for single or recurring use, and the source of the information.
			\item \textbf{Why A is wrong:} Strategy must include more than just collection methods
			\item \textbf{Why C is wrong:} Storage is just one component
			\item \textbf{Why D is wrong:} Deletion is just one component
		\end{itemize}
		
		\textcolor{myred}{\textbf{EXAM TRAP:}} Data strategy requires \textcolor{myblue}{policies, vision, practices, and authorization} - not just technical aspects!
	\end{frame}
	
	\begin{frame}{Q7: Alternative Data Controls - Tutorial}
		\fontsize{6.5}{7.5}\selectfont
		\textbf{QUESTION: Why may alternative data sources require additional controls and testing?}
		
		\vspace{0.1cm}
		\textbf{OPTIONS:}
		\begin{enumerate}
			\item Alternative data is always lower quality than traditional data
			\item Concerns about privacy, data quality, accuracy, and the general nature of unbounded data streams
			\item Alternative data never requires validation
			\item Alternative data is always perfectly clean
		\end{enumerate}
		
		\vspace{0.1cm}
		\textcolor{myblue}{\textbf{ANSWER: B}}
		
		\vspace{0.1cm}
		\textbf{DETAILED EXPLANATION:}
		\begin{itemize}
			\item \textbf{Why B is correct:} The use of alternative data may require additional controls and testing due to concerns about privacy, data quality and accuracy, and other potential risks, as well as the general nature of unbounded data streams.
			\item \textbf{Why A is wrong:} Alternative data is not always lower quality - it just requires more scrutiny
			\item \textbf{Why C is wrong:} Alternative data requires validation like any other data
			\item \textbf{Why D is wrong:} Alternative data can have significant quality issues
		\end{itemize}
		
		\textcolor{myred}{\textbf{EXAM TRAP:}} Alternative data requires \textcolor{myblue}{additional scrutiny} due to privacy, quality, and unbounded nature!
	\end{frame}
	
	\begin{frame}{Q8: Data Quality - Tutorial}
		\fontsize{6.5}{7.5}\selectfont
		\textbf{QUESTION: What should be the primary objective of data governance?}
		
		\vspace{0.1cm}
		\textbf{OPTIONS:}
		\begin{enumerate}
			\item Maximizing the amount of data collected
			\item Complete, relevant, transparent, and accurate data
			\item Minimizing the cost of data storage
			\item Maximizing data sharing with all partners
		\end{enumerate}
		
		\vspace{0.1cm}
		\textcolor{myblue}{\textbf{ANSWER: B}}
		
		\vspace{0.1cm}
		\textbf{DETAILED EXPLANATION:}
		\begin{itemize}
			\item \textbf{Why B is correct:} Complete, relevant, transparent, and accurate data should be the primary objective. This ensures data is fit for purpose and reliable.
			\item \textbf{Why A is wrong:} More data is not always better - quality matters more than quantity
			\item \textbf{Why C is wrong:} Cost is a consideration but not the primary objective
			\item \textbf{Why D is wrong:} Sharing must be controlled and appropriate
		\end{itemize}
		
		\textcolor{myred}{\textbf{EXAM TRAP:}} Data quality means \textcolor{myblue}{complete, relevant, transparent, and accurate} data!
	\end{frame}
	
	\begin{frame}{Q9: Data Provenance - Tutorial}
		\fontsize{6.5}{7.5}\selectfont
		\textbf{QUESTION: What does "data provenance" refer to in the context of AI models?}
		
		\vspace{0.1cm}
		\textbf{OPTIONS:}
		\begin{enumerate}
			\item The age of the data
			\item The sequence of ownership and handling of data, verifying the legal right to use it
			\item The format of the data
			\item The size of the dataset
		\end{enumerate}
		
		\vspace{0.1cm}
		\textcolor{myblue}{\textbf{ANSWER: B}}
		
		\vspace{0.1cm}
		\textbf{DETAILED EXPLANATION:}
		\begin{itemize}
			\item \textbf{Why B is correct:} Provenance refers to the sequence of ownership and handling of items of value. This term now applies to the data used to create models. The use of alternative data has resulted in scrutiny of whether vendors and model builders have a legal right to use this data.
			\item \textbf{Why A is wrong:} Age is different from provenance
			\item \textbf{Why C is wrong:} Format is different from provenance
			\item \textbf{Why D is wrong:} Size is different from provenance
		\end{itemize}
		
		\textcolor{myred}{\textbf{EXAM TRAP:}} Data provenance is about \textcolor{myblue}{legal ownership and handling history} - not technical characteristics!
	\end{frame}
	
	\begin{frame}{Q10: FTC and Improper Data - Tutorial}
		\fontsize{6.5}{7.5}\selectfont
		\textbf{QUESTION: What has the US Federal Trade Commission stated about models built on improperly obtained data?}
		
		\vspace{0.1cm}
		\textbf{OPTIONS:}
		\begin{enumerate}
			\item They can continue to be used
			\item They must be destroyed
			\item They just need to be revalidated
			\item There are no consequences
		\end{enumerate}
		
		\vspace{0.1cm}
		\textcolor{myblue}{\textbf{ANSWER: B}}
		
		\vspace{0.1cm}
		\textbf{DETAILED EXPLANATION:}
		\begin{itemize}
			\item \textbf{Why B is correct:} The US Federal Trade Commission has stated that any model proven to be built on data that was not properly obtained must be destroyed. The first such legal action has already been settled with the destruction of such models.
			\item \textbf{Why A is wrong:} Models on improper data cannot continue to be used
			\item \textbf{Why C is wrong:} Revalidation cannot fix improper data provenance
			\item \textbf{Why D is wrong:} There are significant legal consequences
		\end{itemize}
		
		\textcolor{myred}{\textbf{EXAM TRAP:}} Models built on \textcolor{myblue}{improperly obtained data must be destroyed} - severe consequence!
	\end{frame}
	
	\begin{frame}{Q11: Synthetic Data - Tutorial}
		\fontsize{6.5}{7.5}\selectfont
		\textbf{QUESTION: What is an advantage of using synthetic data in AI model training?}
		
		\vspace{0.1cm}
		\textbf{OPTIONS:}
		\begin{enumerate}
			\item It always perfectly represents real data
			\item It doesn't entail privacy risks since the data does not come from real data subjects
			\item It is always more accurate than real data
			\item It requires no validation
		\end{enumerate}
		
		\vspace{0.1cm}
		\textcolor{myblue}{\textbf{ANSWER: B}}
		
		\vspace{0.1cm}
		\textbf{DETAILED EXPLANATION:}
		\begin{itemize}
			\item \textbf{Why B is correct:} The advantages of synthetic data include that it doesn't entail privacy risks, given that the data does not come from real data subjects. Synthetic data can also be more easily rid of problematic biases by design.
			\item \textbf{Why A is wrong:} Synthetic data may not always perfectly represent real data
			\item \textbf{Why C is wrong:} Synthetic data may be less precise than real data
			\item \textbf{Why D is wrong:} Synthetic data still requires validation
		\end{itemize}
		
		\textcolor{myred}{\textbf{EXAM TRAP:}} Synthetic data offers \textcolor{myblue}{privacy advantages} but may have quality limitations!
	\end{frame}
	
	\begin{frame}{Q12: Data Classification - Tutorial}
		\fontsize{6.5}{7.5}\selectfont
		\textbf{QUESTION: What is critical for good governance regarding data classification?}
		
		\vspace{0.1cm}
		\textbf{OPTIONS:}
		\begin{enumerate}
			\item Collecting as much data as possible
			\item Distinguishing the classification of data being used and ensuring access and use is permissible
			\item Making all data public
			\item Storing all data indefinitely
		\end{enumerate}
		
		\vspace{0.1cm}
		\textcolor{myblue}{\textbf{ANSWER: B}}
		
		\vspace{0.1cm}
		\textbf{DETAILED EXPLANATION:}
		\begin{itemize}
			\item \textbf{Why B is correct:} Data will be either quantitative or qualitative. What is critical for good governance is to distinguish the classification of data being used, and ensure that access and use is permissible. When data is either confidential or personally identifiable, controls must be present, operational, and effective.
			\item \textbf{Why A is wrong:} Collecting more data doesn't help with classification
			\item \textbf{Why C is wrong:} Making data public violates privacy
			\item \textbf{Why D is wrong:} Indefinite storage increases risk
		\end{itemize}
		
		\textcolor{myred}{\textbf{EXAM TRAP:}} Data classification determines \textcolor{myblue}{access and usage permissions}!
	\end{frame}
	
	\begin{frame}{Q13: Metadata Management - Tutorial}
		\fontsize{6.5}{7.5}\selectfont
		\textbf{QUESTION: What does metadata give basic information about?}
		
		\vspace{0.1cm}
		\textbf{OPTIONS:}
		\begin{enumerate}
			\item Only the data's values
			\item Data including file type, time of creation, data author, data source, and file size
			\item Only the data's numerical values
			\item Only the data's column names
		\end{enumerate}
		
		\vspace{0.1cm}
		\textcolor{myblue}{\textbf{ANSWER: B}}
		
		\vspace{0.1cm}
		\textbf{DETAILED EXPLANATION:}
		\begin{itemize}
			\item \textbf{Why B is correct:} Metadata gives basic information about data, including file type, time of creation, data author, data source, file size, and other relevant information.
			\item \textbf{Why A is wrong:} Metadata includes more than just values
			\item \textbf{Why C is wrong:} Metadata includes non-numerical information
			\item \textbf{Why D is wrong:} Metadata includes more than column names
		\end{itemize}
		
		\textcolor{myred}{\textbf{EXAM TRAP:}} Metadata includes \textcolor{myblue}{file type, creation time, author, source, and size} - not just values!
	\end{frame}
	
	\begin{frame}{Q14: Types of Metadata - Tutorial}
		\fontsize{6.5}{7.5}\selectfont
		\textbf{QUESTION: Which is NOT a type of metadata?}
		
		\vspace{0.1cm}
		\textbf{OPTIONS:}
		\begin{enumerate}
			\item Descriptive metadata
			\item Structural metadata
			\item Administrative metadata
			\item Operational metadata
		\end{enumerate}
		
		\vspace{0.1cm}
		\textcolor{myblue}{\textbf{ANSWER: D}}
		
		\vspace{0.1cm}
		\textbf{DETAILED EXPLANATION:}
		\begin{itemize}
			\item \textbf{Why D is correct:} Operational metadata is not one of the standard types. The distinct types of metadata include descriptive, structural, administrative, reference, and statistical metadata.
			\item \textbf{Why A is wrong:} Descriptive IS a type of metadata
			\item \textbf{Why B is wrong:} Structural IS a type of metadata
			\item \textbf{Why C is wrong:} Administrative IS a type of metadata
		\end{itemize}
		
		\textcolor{myred}{\textbf{EXAM TRAP:}} Metadata types include \textcolor{myblue}{descriptive, structural, administrative, reference, and statistical} - not operational!
	\end{frame}
	
	\begin{frame}{Q15: Data Security vs Protection - Tutorial}
		\fontsize{6.5}{7.5}\selectfont
		\textbf{QUESTION: How are data security and data protection related?}
		
		\vspace{0.1cm}
		\textbf{OPTIONS:}
		\begin{enumerate}
			\item They are completely unrelated concepts
			\item Data security is often used interchangeably with data protection and data integrity
			\item Data protection is only about privacy
			\item Data security only applies to cloud storage
		\end{enumerate}
		
		\vspace{0.1cm}
		\textcolor{myblue}{\textbf{ANSWER: B}}
		
		\vspace{0.1cm}
		\textbf{DETAILED EXPLANATION:}
		\begin{itemize}
			\item \textbf{Why B is correct:} Data security is often used interchangeably with data protection and data integrity. The strategic security process, data protection plan, and procedural steps identified to safeguard the availability of data, access to data, and data privacy are all part of data security.
			\item \textbf{Why A is wrong:} They are closely related concepts
			\item \textbf{Why C is wrong:} Data protection includes more than just privacy
			\item \textbf{Why D is wrong:} Data security applies to all data storage
		\end{itemize}
		
		\textcolor{myred}{\textbf{EXAM TRAP:}} Data security and data protection are \textcolor{myblue}{often used interchangeably}!
	\end{frame}
	
	\begin{frame}{Q16: Data Retention Policy - Tutorial}
		\fontsize{6.5}{7.5}\selectfont
		\textbf{QUESTION: What should be in place regarding data retention?}
		
		\vspace{0.1cm}
		\textbf{OPTIONS:}
		\begin{enumerate}
			\item No data retention policy is needed
			\item A data retention policy should be in place and adhered to
			\item Data should be kept forever
			\item Data should be deleted immediately after collection
		\end{enumerate}
		
		\vspace{0.1cm}
		\textcolor{myblue}{\textbf{ANSWER: B}}
		
		\vspace{0.1cm}
		\textbf{DETAILED EXPLANATION:}
		\begin{itemize}
			\item \textbf{Why B is correct:} A data retention policy should be in place and adhered to. This ensures data is kept for appropriate periods and deleted when no longer needed.
			\item \textbf{Why A is wrong:} Retention policies are essential for governance
			\item \textbf{Why C is wrong:} Keeping data forever creates unnecessary risk
			\item \textbf{Why D is wrong:} Immediate deletion may not be practical or legal
		\end{itemize}
		
		\textcolor{myred}{\textbf{EXAM TRAP:}} A data retention policy is \textcolor{myblue}{essential and must be followed}!
	\end{frame}
	
	\begin{frame}{Q17: Data Access - Tutorial}
		\fontsize{6.5}{7.5}\selectfont
		\textbf{QUESTION: What is the most important aspect of data access?}
		
		\vspace{0.1cm}
		\textbf{OPTIONS:}
		\begin{enumerate}
			\item Making all data public
			\item Ensuring data is readily accessible and shareable to those who have requisite permission and need to access it
			\item Restricting all data access
			\item Allowing unlimited access to everyone
		\end{enumerate}
		
		\vspace{0.1cm}
		\textcolor{myblue}{\textbf{ANSWER: B}}
		
		\vspace{0.1cm}
		\textbf{DETAILED EXPLANATION:}
		\begin{itemize}
			\item \textbf{Why B is correct:} Ensuring that data is readily accessible and shareable to those who have the requisite permission and need to access it is the most important aspect of data access.
			\item \textbf{Why A is wrong:} Making data public violates privacy
			\item \textbf{Why C is wrong:} Restricting all access prevents legitimate use
			\item \textbf{Why D is wrong:} Unlimited access violates security principles
		\end{itemize}
		
		\textcolor{myred}{\textbf{EXAM TRAP:}} Data access balances \textcolor{myblue}{availability with permissions and need-to-know}!
	\end{frame}
	
	\begin{frame}{Q18: GLBA - Tutorial}
		\fontsize{6.5}{7.5}\selectfont
		\textbf{QUESTION: What does the Gramm-Leach-Bliley Act (GLBA) require?}
		
		\vspace{0.1cm}
		\textbf{OPTIONS:}
		\begin{enumerate}
			\item Nothing regarding data privacy
			\item Financial institutions to provide customers with information about privacy practices, opt-out rights, and implement security safeguards
			\item Only healthcare providers to protect data
			\item Only government agencies to protect data
		\end{enumerate}
		
		\vspace{0.1cm}
		\textcolor{myblue}{\textbf{ANSWER: B}}
		
		\vspace{0.1cm}
		\textbf{DETAILED EXPLANATION:}
		\begin{itemize}
			\item \textbf{Why B is correct:} The Gramm-Leach-Bliley Act (GLBA) requires financial institutions to provide customers with information about the institutions' privacy practices and about their opt-out rights, and to implement security safeguards for customer information.
			\item \textbf{Why A is wrong:} GLBA has significant privacy requirements
			\item \textbf{Why C is wrong:} GLBA applies to financial institutions, not healthcare
			\item \textbf{Why D is wrong:} GLBA applies to financial institutions
		\end{itemize}
		
		\textcolor{myred}{\textbf{EXAM TRAP:}} GLBA requires \textcolor{myblue}{privacy notices, opt-out rights, and security safeguards} for financial institutions!
	\end{frame}
	
	\begin{frame}{Q19: Data Compliance Principles - Tutorial}
		\fontsize{6.5}{7.5}\selectfont
		\textbf{QUESTION: What are the three basic principles of data compliance regulations?}
		
		\vspace{0.1cm}
		\textbf{OPTIONS:}
		\begin{enumerate}
			\item Collect all data, store forever, share widely
			\item Obtaining consent, minimizing data held, ensuring rights of data subjects
			\item Only collecting data, only storing data, only deleting data
			\item Ignoring consent, maximizing data, disregarding rights
		\end{enumerate}
		
		\vspace{0.1cm}
		\textcolor{myblue}{\textbf{ANSWER: B}}
		
		\vspace{0.1cm}
		\textbf{DETAILED EXPLANATION:}
		\begin{itemize}
			\item \textbf{Why B is correct:} Although there are many rules within current regulations, the majority can be boiled down to three basic principles: obtaining consent, minimizing the amount of data you hold, and ensuring the rights of data subjects.
			\item \textbf{Why A is wrong:} This violates all privacy principles
			\item \textbf{Why C is wrong:} These are technical activities, not principles
			\item \textbf{Why D is wrong:} This is the opposite of compliance
		\end{itemize}
		
		\textcolor{myred}{\textbf{EXAM TRAP:}} Data compliance principles are \textcolor{myblue}{consent, minimization, and data subject rights}!
	\end{frame}
	
	\begin{frame}{Q20: Board Data Governance Role - Tutorial}
		\fontsize{6.5}{7.5}\selectfont
		\textbf{QUESTION: What is the Board's role in data governance?}
		
		\vspace{0.1cm}
		\textbf{OPTIONS:}
		\begin{enumerate}
			\item No role - only IT department handles data governance
			\item Overseeing the data governance framework, ensuring alignment with strategic objectives, risk management policies, and compliance requirements
			\item Only approving data storage purchases
			\item Only reviewing data security once per year
		\end{enumerate}
		
		\vspace{0.1cm}
		\textcolor{myblue}{\textbf{ANSWER: B}}
		
		\vspace{0.1cm}
		\textbf{DETAILED EXPLANATION:}
		\begin{itemize}
			\item \textbf{Why B is correct:} The Board plays a crucial role in overseeing a firm's data governance framework, and ensuring that the framework aligns with the organization's strategic objectives, risk management policies, and compliance requirements.
			\item \textbf{Why A is wrong:} The Board has significant oversight responsibility
			\item \textbf{Why C is wrong:} Data governance involves more than just storage
			\item \textbf{Why D is wrong:} Oversight is continuous, not just annual
		\end{itemize}
		
		\textcolor{myred}{\textbf{EXAM TRAP:}} The Board sets the \textcolor{myblue}{"tone at the top"} for data governance culture!
	\end{frame}
	
	% ============================================================
	% SECTION 3: MODEL GOVERNANCE - QUESTIONS 21-35
	% ============================================================
	
	\begin{frame}{Q21: Model Definition - Tutorial}
		\fontsize{6.5}{7.5}\selectfont
		\textbf{QUESTION: How does SR 11-7 define a model?}
		
		\vspace{0.1cm}
		\textbf{OPTIONS:}
		\begin{enumerate}
			\item Any spreadsheet used in banking
			\item A quantitative method, system, or approach that applies statistical, economic, financial, or mathematical theories, techniques, and assumptions to process input data into quantitative estimates
			\item Only neural networks and deep learning systems
			\item Any software program
		\end{enumerate}
		
		\vspace{0.1cm}
		\textcolor{myblue}{\textbf{ANSWER: B}}
		
		\vspace{0.1cm}
		\textbf{DETAILED EXPLANATION:}
		\begin{itemize}
			\item \textbf{Why B is correct:} SR 11-7 defines a model as "a quantitative method, system, or approach that applies statistical, economic, financial, or mathematical theories, techniques, and assumptions to process input data into quantitative estimates."
			\item \textbf{Why A is wrong:} Not all spreadsheets qualify as models
			\item \textbf{Why C is wrong:} Models include many types beyond neural networks
			\item \textbf{Why D is wrong:} Not all software programs are models
		\end{itemize}
		
		\textcolor{myred}{\textbf{EXAM TRAP:}} SR 11-7 defines a model as a \textcolor{myblue}{quantitative method using theories and assumptions to process data into estimates}!
	\end{frame}
	
	\begin{frame}{Q22: Risk Modeling Elements - Tutorial}
		\fontsize{6.5}{7.5}\selectfont
		\textbf{QUESTION: What are the three main elements of quantitative risk modeling?}
		
		\vspace{0.1cm}
		\textbf{OPTIONS:}
		\begin{enumerate}
			\item Input data, model parameters, output reports
			\item Quantity of interest, potential future scenarios, risk measure
			\item Training data, validation data, test data
			\item Development, validation, deployment
		\end{enumerate}
		
		\vspace{0.1cm}
		\textcolor{myblue}{\textbf{ANSWER: B}}
		
		\vspace{0.1cm}
		\textbf{DETAILED EXPLANATION:}
		\begin{itemize}
			\item \textbf{Why B is correct:} Quantitative risk modeling encompasses three main elements: (1) Quantity of interest - a numerical object whose future value is uncertain, (2) Potential future scenarios - possible values of the quantity of interest with assigned weights, and (3) Risk measure - summarizes essential information derived from analyzing the potential future scenarios.
			\item \textbf{Why A is wrong:} These are technical components, not the conceptual elements
			\item \textbf{Why C is wrong:} These are data splits for model development
			\item \textbf{Why D is wrong:} These are lifecycle phases
		\end{itemize}
		
		\textcolor{myred}{\textbf{EXAM TRAP:}} Risk modeling has \textcolor{myblue}{three conceptual elements}: quantity of interest, scenarios, and risk measure!
	\end{frame}
	
	\begin{frame}{Q23: Model Governance Policies - Tutorial}
		\fontsize{6.5}{7.5}\selectfont
		\textbf{QUESTION: What should model governance policies address?}
		
		\vspace{0.1cm}
		\textbf{OPTIONS:}
		\begin{enumerate}
			\item Only model development
			\item Defining what constitutes a model, model risk definition, risk tier identification, validation definition, and documentation standards
			\item Only model validation
			\item Only model deployment
		\end{enumerate}
		
		\vspace{0.1cm}
		\textcolor{myblue}{\textbf{ANSWER: B}}
		
		\vspace{0.1cm}
		\textbf{DETAILED EXPLANATION:}
		\begin{itemize}
			\item \textbf{Why B is correct:} Model governance policies should address: defining what constitutes a model, defining model risk, identification of risk tier, establishing validation definition, specifying information for risk assessment, documenting risk appetite, documentation standards, performance monitoring standards, and risk reporting.
			\item \textbf{Why A is wrong:} Governance covers more than just development
			\item \textbf{Why C is wrong:} Governance covers more than just validation
			\item \textbf{Why D is wrong:} Governance covers more than just deployment
		\end{itemize}
		
		\textcolor{myred}{\textbf{EXAM TRAP:}} Model governance policies are \textcolor{myblue}{comprehensive across the entire lifecycle}!
	\end{frame}
	
	\begin{frame}{Q24: Model Documentation Standards - Tutorial}
		\fontsize{6.5}{7.5}\selectfont
		\textbf{QUESTION: According to SR 11-7, how detailed should model documentation be?}
		
		\vspace{0.1cm}
		\textbf{OPTIONS:}
		\begin{enumerate}
			\item Only a one-page summary
			\item Sufficiently detailed to allow parties unfamiliar with a model to understand how it operates, its limitations, and key assumptions
			\item Only the code is needed
			\item Documentation is not required
		\end{enumerate}
		
		\vspace{0.1cm}
		\textcolor{myblue}{\textbf{ANSWER: B}}
		
		\vspace{0.1cm}
		\textbf{DETAILED EXPLANATION:}
		\begin{itemize}
			\item \textbf{Why B is correct:} SR 11-7 states: "Strong governance also includes documentation of model development and validation that is sufficiently detailed to allow parties unfamiliar with a model to understand how the model operates, as well as its limitations and key assumptions."
			\item \textbf{Why A is wrong:} A one-page summary is insufficient
			\item \textbf{Why C is wrong:} Documentation must include more than just code
			\item \textbf{Why D is wrong:} Documentation is required for strong governance
		\end{itemize}
		
		\textcolor{myred}{\textbf{EXAM TRAP:}} Documentation must be \textcolor{myblue}{detailed enough for unfamiliar parties to understand the model}!
	\end{frame}
	
	\begin{frame}{Q25: Model Inventory - Tutorial}
		\fontsize{6.5}{7.5}\selectfont
		\textbf{QUESTION: What is the purpose of a model inventory?}
		
		\vspace{0.1cm}
		\textbf{OPTIONS:}
		\begin{enumerate}
			\item To make models harder to find
			\item To enhance transparency regarding the number of models in use and track their usage, changes, and approvals
			\item To hide models from regulators
			\item To reduce the number of models
		\end{enumerate}
		
		\vspace{0.1cm}
		\textcolor{myblue}{\textbf{ANSWER: B}}
		
		\vspace{0.1cm}
		\textbf{DETAILED EXPLANATION:}
		\begin{itemize}
			\item \textbf{Why B is correct:} To facilitate model governance, institutions should maintain a model inventory, which enhances transparency regarding the number of models in use and tracks their usage, changes, and approvals.
			\item \textbf{Why A is wrong:} The purpose is to make models more transparent
			\item \textbf{Why C is wrong:} The purpose is to enhance transparency for regulators
			\item \textbf{Why D is wrong:} The purpose is tracking, not reducing
		\end{itemize}
		
		\textcolor{myred}{\textbf{EXAM TRAP:}} Model inventory enhances \textcolor{myblue}{transparency and tracking} of all models!
	\end{frame}
	
	\begin{frame}{Q26: Model Inventory Content - Tutorial}
		\fontsize{6.5}{7.5}\selectfont
		\textbf{QUESTION: What should a comprehensive model inventory include?}
		
		\vspace{0.1cm}
		\textbf{OPTIONS:}
		\begin{enumerate}
			\item Only the model name
			\item Model owner, purpose, methodology, classification, development source, assumptions, limitations, documentation locations, and validation history
			\item Only the model code
			\item Only the model's last validation date
		\end{enumerate}
		
		\vspace{0.1cm}
		\textcolor{myblue}{\textbf{ANSWER: B}}
		
		\vspace{0.1cm}
		\textbf{DETAILED EXPLANATION:}
		\begin{itemize}
			\item \textbf{Why B is correct:} A comprehensive model inventory should include: model owner, user, and developer information; purpose and description; methodology; classification; development source; use and frequency; restrictions; assumptions; limitations; documentation locations; code access; interdependencies; update history; validation findings; performance monitoring; and risk tier.
			\item \textbf{Why A is wrong:} Much more than the name is needed
			\item \textbf{Why C is wrong:} Code is just one component
			\item \textbf{Why D is wrong:} Much more than the last validation date is needed
		\end{itemize}
		
		\textcolor{myred}{\textbf{EXAM TRAP:}} Model inventory is \textcolor{myblue}{comprehensive} - not just basic information!
	\end{frame}
	
	\begin{frame}{Q27: Model Landscape - Tutorial}
		\fontsize{6.5}{7.5}\selectfont
		\textbf{QUESTION: What does a model landscape illustrate?}
		
		\vspace{0.1cm}
		\textbf{OPTIONS:}
		\begin{enumerate}
			\item Only the number of models
			\item The interactions among models, emphasizing crucial models and their impact on key business decisions
			\item Only model validation dates
			\item Only model owners
		\end{enumerate}
		
		\vspace{0.1cm}
		\textcolor{myblue}{\textbf{ANSWER: B}}
		
		\vspace{0.1cm}
		\textbf{DETAILED EXPLANATION:}
		\begin{itemize}
			\item \textbf{Why B is correct:} A model landscape is a graphical or list-based representation that illustrates the interactions among models, emphasizing crucial models and their impact on key business decisions. While maintaining clarity, model landscapes should avoid overwhelming users with excessive information.
			\item \textbf{Why A is wrong:} It shows interactions, not just counts
			\item \textbf{Why C is wrong:} Validation dates are not the main focus
			\item \textbf{Why D is wrong:} Owners are just one element
		\end{itemize}
		
		\textcolor{myred}{\textbf{EXAM TRAP:}} Model landscape shows \textcolor{myblue}{interactions and interdependencies} among models!
	\end{frame}
	
	\begin{frame}{Q28: AI/ML Registration Considerations - Tutorial}
		\fontsize{6.5}{7.5}\selectfont
		\textbf{QUESTION: What new aspects must be considered when registering AI/ML models?}
		
		\vspace{0.1cm}
		\textbf{OPTIONS:}
		\begin{enumerate}
			\item Only the model name
			\item Complexity of methodology, data usage, output parameters, and recalibration needs
			\item Only the model's accuracy
			\item Only the model's speed
		\end{enumerate}
		
		\vspace{0.1cm}
		\textcolor{myblue}{\textbf{ANSWER: B}}
		
		\vspace{0.1cm}
		\textbf{DETAILED EXPLANATION:}
		\begin{itemize}
			\item \textbf{Why B is correct:} When registering AI/ML models, more weight must be put on: complexity of methodology and design, data usage (volume, complexity, quality), output parameters (supervised vs. unsupervised), and model recalibration needs (static vs. continuous).
			\item \textbf{Why A is wrong:} Much more than the name is needed
			\item \textbf{Why C is wrong:} Accuracy is just one aspect
			\item \textbf{Why D is wrong:} Speed is just one aspect
		\end{itemize}
		
		\textcolor{myred}{\textbf{EXAM TRAP:}} AI/ML registration requires \textcolor{myblue}{additional considerations beyond traditional models}!
	\end{frame}
	
	\begin{frame}{Q29: Model Risk Ranking Factors - Tutorial}
		\fontsize{6.5}{7.5}\selectfont
		\textbf{QUESTION: What are the key factors in model risk ranking?}
		
		\vspace{0.1cm}
		\textbf{OPTIONS:}
		\begin{enumerate}
			\item Only the model's age
			\item Materiality, complexity, and exposure
			\item Only the model's accuracy
			\item Only the number of users
		\end{enumerate}
		
		\vspace{0.1cm}
		\textcolor{myblue}{\textbf{ANSWER: B}}
		
		\vspace{0.1cm}
		\textbf{DETAILED EXPLANATION:}
		\begin{itemize}
			\item \textbf{Why B is correct:} Model risk ranking factors include materiality (economic, operational, reputational consequences), complexity (which may elevate risk ranking even if materiality is low), and exposure (if a simple model affects published financial reports, risk ranking may be elevated).
			\item \textbf{Why A is wrong:} Age is not a primary factor
			\item \textbf{Why C is wrong:} Accuracy is just one aspect
			\item \textbf{Why D is wrong:} Number of users is just one aspect
		\end{itemize}
		
		\textcolor{myred}{\textbf{EXAM TRAP:}} Risk ranking considers \textcolor{myblue}{materiality, complexity, and exposure}!
	\end{frame}
	
	\begin{frame}{Q30: Minimum Standards - Tutorial}
		\fontsize{6.5}{7.5}\selectfont
		\textbf{QUESTION: Should minimum standards of model risk assessment apply to all models?}
		
		\vspace{0.1cm}
		\textbf{OPTIONS:}
		\begin{enumerate}
			\item No, only high-risk models need assessment
			\item Yes, minimum standards should apply to all models regardless of materiality
			\item Only models used for regulatory reporting need assessment
			\item Only vendor models need assessment
		\end{enumerate}
		
		\vspace{0.1cm}
		\textcolor{myblue}{\textbf{ANSWER: B}}
		
		\vspace{0.1cm}
		\textbf{DETAILED EXPLANATION:}
		\begin{itemize}
			\item \textbf{Why B is correct:} While materiality, complexity, and exposure assessments are crucial, minimum standards of model risk assessment should be applied to all models regardless of their materiality, as the absence of such assessments can introduce its own model risk.
			\item \textbf{Why A is wrong:} All models need assessment
			\item \textbf{Why C is wrong:} All models need assessment, not just regulatory ones
			\item \textbf{Why D is wrong:} Both in-house and vendor models need assessment
		\end{itemize}
		
		\textcolor{myred}{\textbf{EXAM TRAP:}} Minimum standards apply to \textcolor{myblue}{all models regardless of materiality}!
	\end{frame}
	
	\begin{frame}{Q31: First Line of Defense - Tutorial}
		\fontsize{6.5}{7.5}\selectfont
		\textbf{QUESTION: Who comprises the First Line of Defense in model risk management?}
		
		\vspace{0.1cm}
		\textbf{OPTIONS:}
		\begin{enumerate}
			\item Internal audit
			\item Business and model owners/developers
			\item Regulators
			\item External consultants
		\end{enumerate}
		
		\vspace{0.1cm}
		\textcolor{myblue}{\textbf{ANSWER: B}}
		
		\vspace{0.1cm}
		\textbf{DETAILED EXPLANATION:}
		\begin{itemize}
			\item \textbf{Why B is correct:} The First Line of Defense comprises business and model owners/developers. They own the model and its performance, are responsible for designing and implementing models, operate performance dashboards, own data quality for model inputs, and are the first to react to threshold breaches.
			\item \textbf{Why A is wrong:} Internal audit is the Third Line of Defense
			\item \textbf{Why C is wrong:} Regulators are external oversight
			\item \textbf{Why D is wrong:} External consultants are not part of the three lines
		\end{itemize}
		
		\textcolor{myred}{\textbf{EXAM TRAP:}} First Line = \textcolor{myblue}{business and model owners/developers} who own the model!
	\end{frame}
	
	\begin{frame}{Q32: Second Line of Defense - Tutorial}
		\fontsize{6.5}{7.5}\selectfont
		\textbf{QUESTION: Who comprises the Second Line of Defense in model risk management?}
		
		\vspace{0.1cm}
		\textbf{OPTIONS:}
		\begin{enumerate}
			\item Business and model owners
			\item Independent risk management/model validation
			\item Internal audit
			\item Model users
		\end{enumerate}
		
		\vspace{0.1cm}
		\textcolor{myblue}{\textbf{ANSWER: B}}
		
		\vspace{0.1cm}
		\textbf{DETAILED EXPLANATION:}
		\begin{itemize}
			\item \textbf{Why B is correct:} The Second Line of Defense is independent risk management/model validation. They set model risk management policy, conduct independent model validations, challenge the First Line's performance assessments, maintain the enterprise-wide model inventory, and report material model performance issues to senior management.
			\item \textbf{Why A is wrong:} This is the First Line
			\item \textbf{Why C is wrong:} This is the Third Line
			\item \textbf{Why D is wrong:} Model users are in the First Line
		\end{itemize}
		
		\textcolor{myred}{\textbf{EXAM TRAP:}} Second Line provides \textcolor{myblue}{independent challenge and oversight} of models!
	\end{frame}
	
	\begin{frame}{Q33: Third Line of Defense - Tutorial}
		\fontsize{6.5}{7.5}\selectfont
		\textbf{QUESTION: Who comprises the Third Line of Defense in model risk management?}
		
		\vspace{0.1cm}
		\textbf{OPTIONS:}
		\begin{enumerate}
			\item Business and model owners
			\item Independent risk management
			\item Internal audit
			\item Model developers
		\end{enumerate}
		
		\vspace{0.1cm}
		\textcolor{myblue}{\textbf{ANSWER: C}}
		
		\vspace{0.1cm}
		\textbf{DETAILED EXPLANATION:}
		\begin{itemize}
			\item \textbf{Why C is correct:} The Third Line of Defense is internal audit. They provide independent assurance to the board that the model risk management framework, including performance monitoring and validation processes across the First and Second Lines, is well-designed and operating effectively.
			\item \textbf{Why A is wrong:} This is the First Line
			\item \textbf{Why B is wrong:} This is the Second Line
			\item \textbf{Why D is wrong:} Model developers are in the First Line
		\end{itemize}
		
		\textcolor{myred}{\textbf{EXAM TRAP:}} Third Line = \textcolor{myblue}{internal audit} providing independent assurance!
	\end{frame}
	
	\begin{frame}{Q34: Model Owner Role - Tutorial}
		\fontsize{6.5}{7.5}\selectfont
		\textbf{QUESTION: What is the role of a model owner?}
		
		\vspace{0.1cm}
		\textbf{OPTIONS:}
		\begin{enumerate}
			\item Only developing the model
			\item Overseeing proper development, validation, approval, findings remediation, implementation, use, ongoing performance monitoring, and change management
			\item Only using the model
			\item Only validating the model
		\end{enumerate}
		
		\vspace{0.1cm}
		\textcolor{myblue}{\textbf{ANSWER: B}}
		
		\vspace{0.1cm}
		\textbf{DETAILED EXPLANATION:}
		\begin{itemize}
			\item \textbf{Why B is correct:} A model owner oversees a number of models, ensuring their proper development, validation, approval, findings remediation, implementation, use, ongoing performance monitoring and reporting, and change management.
			\item \textbf{Why A is wrong:} Ownership includes more than development
			\item \textbf{Why C is wrong:} Ownership includes more than usage
			\item \textbf{Why D is wrong:} Validation is a separate function
		\end{itemize}
		
		\textcolor{myred}{\textbf{EXAM TRAP:}} Model owners oversee \textcolor{myblue}{the entire model lifecycle} - not just one aspect!
	\end{frame}
	
	\begin{frame}{Q35: Model Validator Role - Tutorial}
		\fontsize{6.5}{7.5}\selectfont
		\textbf{QUESTION: What is the role of model validators?}
		
		\vspace{0.1cm}
		\textbf{OPTIONS:}
		\begin{enumerate}
			\item Developing models
			\item Evaluating models and their performance, providing independence and objectivity
			\item Only using models
			\item Only documenting models
		\end{enumerate}
		
		\vspace{0.1cm}
		\textcolor{myblue}{\textbf{ANSWER: B}}
		
		\vspace{0.1cm}
		\textbf{DETAILED EXPLANATION:}
		\begin{itemize}
			\item \textbf{Why B is correct:} Model validators play a critical role in evaluating models and their performance. Validators provide independence and objectivity within the validation process, and can add value in terms of reviewing assumptions, methodologies, and so on.
			\item \textbf{Why A is wrong:} Development is a different role
			\item \textbf{Why C is wrong:} Usage is a different role
			\item \textbf{Why D is wrong:} Documentation is part of development
		\end{itemize}
		
		\textcolor{myred}{\textbf{EXAM TRAP:}} Validators provide \textcolor{myblue}{independent and objective evaluation} of models!
	\end{frame}
	
	% ============================================================
	% SECTION 4: MODEL DEVELOPMENT AND TESTING - QUESTIONS 36-50
	% ============================================================
	
	\begin{frame}{Q36: Model Development Steps - Tutorial}
		\fontsize{6.5}{7.5}\selectfont
		\textbf{QUESTION: What is the first step in model development?}
		
		\vspace{0.1cm}
		\textbf{OPTIONS:}
		\begin{enumerate}
			\item Data collection
			\item Define objectives and scope
			\item Model training
			\item Feature engineering
		\end{enumerate}
		
		\vspace{0.1cm}
		\textcolor{myblue}{\textbf{ANSWER: B}}
		
		\vspace{0.1cm}
		\textbf{DETAILED EXPLANATION:}
		\begin{itemize}
			\item \textbf{Why B is correct:} The first step in model development is to define the objectives and scope. Clearly define the objectives and scope of the model, determine the specific problem or risk being addressed, the applicable products, the required data, desired outcomes, and the target audience for the model's outputs.
			\item \textbf{Why A is wrong:} Data collection comes after defining objectives
			\item \textbf{Why C is wrong:} Training comes much later in the process
			\item \textbf{Why D is wrong:} Feature engineering comes after data analysis
		\end{itemize}
		
		\textcolor{myred}{\textbf{EXAM TRAP:}} Model development starts with \textcolor{myblue}{defining objectives and scope} - not data collection!
	\end{frame}
	
	\begin{frame}{Q37: Data Leakage - Tutorial}
		\fontsize{6.5}{7.5}\selectfont
		\textbf{QUESTION: What should be avoided during model training?}
		
		\vspace{0.1cm}
		\textbf{OPTIONS:}
		\begin{enumerate}
			\item Using cross-validation
			\item Data leakage (using future or concurrent data that wouldn't have been observable at the point of training)
			\item Testing the model
			\item Documenting the model
		\end{enumerate}
		
		\vspace{0.1cm}
		\textcolor{myblue}{\textbf{ANSWER: B}}
		
		\vspace{0.1cm}
		\textbf{DETAILED EXPLANATION:}
		\begin{itemize}
			\item \textbf{Why B is correct:} Special emphasis should be put on training with data leakage (using future or concurrent data that wouldn't have been observable at the point of the training data). Use accepted techniques such as cross-validation.
			\item \textbf{Why A is wrong:} Cross-validation is recommended, not avoided
			\item \textbf{Why C is wrong:} Testing is essential, not avoided
			\item \textbf{Why D is wrong:} Documentation is essential, not avoided
		\end{itemize}
		
		\textcolor{myred}{\textbf{EXAM TRAP:}} Data leakage occurs when \textcolor{myblue}{future data is used in training} - must be avoided!
	\end{frame}
	
	\begin{frame}{Q38: Testing Responsibility - Tutorial}
		\fontsize{6.5}{7.5}\selectfont
		\textbf{QUESTION: Who is responsible for testing models?}
		
		\vspace{0.1cm}
		\textbf{OPTIONS:}
		\begin{enumerate}
			\item Only the model developer
			\item Programmers, dedicated testers, and external testers with fresh perspectives
			\item Only external consultants
			\item Only the model owner
		\end{enumerate}
		
		\vspace{0.1cm}
		\textcolor{myblue}{\textbf{ANSWER: B}}
		
		\vspace{0.1cm}
		\textbf{DETAILED EXPLANATION:}
		\begin{itemize}
			\item \textbf{Why B is correct:} Programmers bear the responsibility of debugging and resolving known bugs. Testers, who possess expertise and qualifications in software testing, are entrusted with discovering remaining significant bugs. It is common to rely partially on external testers to offer a fresh perspective.
			\item \textbf{Why A is wrong:} Multiple parties are involved in testing
			\item \textbf{Why C is wrong:} Internal testers are also needed
			\item \textbf{Why D is wrong:} The owner oversees but doesn't do all testing
		\end{itemize}
		
		\textcolor{myred}{\textbf{EXAM TRAP:}} Testing involves \textcolor{myblue}{programmers, testers, and external perspectives}!
	\end{frame}
	
	\begin{frame}{Q39: Unit Tests - Tutorial}
		\fontsize{6.5}{7.5}\selectfont
		\textbf{QUESTION: What do unit tests validate?}
		
		\vspace{0.1cm}
		\textbf{OPTIONS:}
		\begin{enumerate}
			\item The entire system
			\item That each piece of code performs as designed
			\item The interfaces between components
			\item User acceptance
		\end{enumerate}
		
		\vspace{0.1cm}
		\textcolor{myblue}{\textbf{ANSWER: B}}
		
		\vspace{0.1cm}
		\textbf{DETAILED EXPLANATION:}
		\begin{itemize}
			\item \textbf{Why B is correct:} Unit tests validate that each piece of the code performs as designed. They test individual components at the smallest level.
			\item \textbf{Why A is wrong:} System tests validate the entire system
			\item \textbf{Why C is wrong:} Integration tests validate interfaces
			\item \textbf{Why D is wrong:} Acceptance tests validate user requirements
		\end{itemize}
		
		\textcolor{myred}{\textbf{EXAM TRAP:}} Unit tests check \textcolor{myblue}{individual code pieces} - foundation of testing!
	\end{frame}
	
	\begin{frame}{Q40: Integration Tests - Tutorial}
		\fontsize{6.5}{7.5}\selectfont
		\textbf{QUESTION: What do integration tests validate?}
		
		\vspace{0.1cm}
		\textbf{OPTIONS:}
		\begin{enumerate}
			\item Individual code pieces
			\item The interfaces between components
			\item The completely integrated system
			\item System performance
		\end{enumerate}
		
		\vspace{0.1cm}
		\textcolor{myblue}{\textbf{ANSWER: B}}
		
		\vspace{0.1cm}
		\textbf{DETAILED EXPLANATION:}
		\begin{itemize}
			\item \textbf{Why B is correct:} Integration tests verify the interfaces between components (e.g., between the QRM's core and its input processing components).
			\item \textbf{Why A is wrong:} Unit tests validate individual pieces
			\item \textbf{Why C is wrong:} System tests validate the complete system
			\item \textbf{Why D is wrong:} Performance tests assess speed and stability
		\end{itemize}
		
		\textcolor{myred}{\textbf{EXAM TRAP:}} Integration tests check \textcolor{myblue}{interfaces between components}!
	\end{frame}
	
	\begin{frame}{Q41: Adversarial Testing - Tutorial}
		\fontsize{6.5}{7.5}\selectfont
		\textbf{QUESTION: What is adversarial testing (red team testing)?}
		
		\vspace{0.1cm}
		\textbf{OPTIONS:}
		\begin{enumerate}
			\item Testing for normal user behavior
			\item Evaluating the response of AI models when subjected to extreme scenarios or edge cases, including prompt injection and unexpected queries
			\item Testing for system speed
			\item Testing for user acceptance
		\end{enumerate}
		
		\vspace{0.1cm}
		\textcolor{myblue}{\textbf{ANSWER: B}}
		
		\vspace{0.1cm}
		\textbf{DETAILED EXPLANATION:}
		\begin{itemize}
			\item \textbf{Why B is correct:} Adversarial testing is performed by "red team" to evaluate the response of AI models when subjected to extreme scenarios or edge cases. It involves tricking the system into making an error by trying out prompt injection, doing a number of chained prompts, or unexpected queries.
			\item \textbf{Why A is wrong:} It's for extreme scenarios, not normal behavior
			\item \textbf{Why C is wrong:} Speed testing is separate
			\item \textbf{Why D is wrong:} User acceptance testing is separate
		\end{itemize}
		
		\textcolor{myred}{\textbf{EXAM TRAP:}} Adversarial testing involves \textcolor{myblue}{extreme scenarios and attempts to trick the system}!
	\end{frame}
	
	\begin{frame}{Q42: White-Box Testing - Tutorial}
		\fontsize{6.5}{7.5}\selectfont
		\textbf{QUESTION: What is white-box testing?}
		
		\vspace{0.1cm}
		\textbf{OPTIONS:}
		\begin{enumerate}
			\item Testing without knowledge of internal implementation
			\item Testing with access to internal data structures, algorithms, and actual code
			\item Testing from a user's perspective
			\item Testing only the user interface
		\end{enumerate}
		
		\vspace{0.1cm}
		\textcolor{myblue}{\textbf{ANSWER: B}}
		
		\vspace{0.1cm}
		\textbf{DETAILED EXPLANATION:}
		\begin{itemize}
			\item \textbf{Why B is correct:} White-box testing involves testers having access to internal data structures, algorithms, and the actual code. It entails testing at unit level from a developer's perspective and may include line-by-line proofreading of the code.
			\item \textbf{Why A is wrong:} This describes black-box testing
			\item \textbf{Why C is wrong:} This describes black-box testing
			\item \textbf{Why D is wrong:} UI testing is just one aspect
		\end{itemize}
		
		\textcolor{myred}{\textbf{EXAM TRAP:}} White-box testing has \textcolor{myblue}{access to internal code and structures}!
	\end{frame}
	
	\begin{frame}{Q43: Black-Box Testing - Tutorial}
		\fontsize{6.5}{7.5}\selectfont
		\textbf{QUESTION: What is black-box testing?}
		
		\vspace{0.1cm}
		\textbf{OPTIONS:}
		\begin{enumerate}
			\item Testing with access to internal code
			\item Testing from a user's perspective without knowledge of internal implementation
			\item Testing only the code
			\item Testing at the unit level
		\end{enumerate}
		
		\vspace{0.1cm}
		\textcolor{myblue}{\textbf{ANSWER: B}}
		
		\vspace{0.1cm}
		\textbf{DETAILED EXPLANATION:}
		\begin{itemize}
			\item \textbf{Why B is correct:} Black-box testing treats the software as a closed box, from a user's perspective, without any knowledge of its internal implementation.
			\item \textbf{Why A is wrong:} This describes white-box testing
			\item \textbf{Why C is wrong:} This is too narrow
			\item \textbf{Why D is wrong:} Unit testing is a type of white-box testing
		\end{itemize}
		
		\textcolor{myred}{\textbf{EXAM TRAP:}} Black-box testing has \textcolor{myblue}{no knowledge of internal implementation} - user perspective!
	\end{frame}
	
	\begin{frame}{Q44: Gray-Box Testing - Tutorial}
		\fontsize{6.5}{7.5}\selectfont
		\textbf{QUESTION: What is gray-box testing?}
		
		\vspace{0.1cm}
		\textbf{OPTIONS:}
		\begin{enumerate}
			\item Testing with full access to code
			\item Testing with no knowledge of implementation
			\item Testing informed by knowledge of functional specifications and architecture but without accessing the code itself
			\item Testing only the user interface
		\end{enumerate}
		
		\vspace{0.1cm}
		\textcolor{myblue}{\textbf{ANSWER: C}}
		
		\vspace{0.1cm}
		\textbf{DETAILED EXPLANATION:}
		\begin{itemize}
			\item \textbf{Why C is correct:} Gray box testing is somewhere between white-box and black-box test, as it is informed by knowledge about functional specifications and the architecture of the underlying code. However, it does not access the code itself.
			\item \textbf{Why A is wrong:} This describes white-box testing
			\item \textbf{Why B is wrong:} This describes black-box testing
			\item \textbf{Why D is wrong:} UI testing is just one aspect
		\end{itemize}
		
		\textcolor{myred}{\textbf{EXAM TRAP:}} Gray-box testing has \textcolor{myblue}{knowledge of architecture but not code access}!
	\end{frame}
	
	\begin{frame}{Q45: Use Test - Tutorial}
		\fontsize{6.5}{7.5}\selectfont
		\textbf{QUESTION: What does the use test examine?}
		
		\vspace{0.1cm}
		\textbf{OPTIONS:}
		\begin{enumerate}
			\item Only the model's mathematical accuracy
			\item The model within its context, considering human interaction, actual usage, acceptance, interpretation of results, and application of those results
			\item Only the model's code
			\item Only the model's performance metrics
		\end{enumerate}
		
		\vspace{0.1cm}
		\textcolor{myblue}{\textbf{ANSWER: B}}
		
		\vspace{0.1cm}
		\textbf{DETAILED EXPLANATION:}
		\begin{itemize}
			\item \textbf{Why B is correct:} The use test examines the model within its context, considering human interaction, actual usage, acceptance of the model, interpretation of results, and the application of those results. It is qualitative in nature and focuses on people rather than numbers.
			\item \textbf{Why A is wrong:} Accuracy is just one aspect
			\item \textbf{Why C is wrong:} Code review is separate
			\item \textbf{Why D is wrong:} Performance metrics are separate
		\end{itemize}
		
		\textcolor{myred}{\textbf{EXAM TRAP:}} Use test focuses on \textcolor{myblue}{how people interact with and use the model} - qualitative assessment!
	\end{frame}
	
	\begin{frame}{Q46: Use Test Reporting - Tutorial}
		\fontsize{6.5}{7.5}\selectfont
		\textbf{QUESTION: How are use test results typically reported?}
		
		\vspace{0.1cm}
		\textbf{OPTIONS:}
		\begin{enumerate}
			\item As detailed technical spreadsheets
			\item Presented to senior management rather than documented as technical reports or spreadsheets
			\item Only to the model development team
			\item They are never documented
		\end{enumerate}
		
		\vspace{0.1cm}
		\textcolor{myblue}{\textbf{ANSWER: B}}
		
		\vspace{0.1cm}
		\textbf{DETAILED EXPLANATION:}
		\begin{itemize}
			\item \textbf{Why B is correct:} The results of the use test are typically presented to senior management rather than documented as technical reports or spreadsheets. Consequently, it is often easier to communicate the results of a use test, minimizing the risk of it being regarded as an arduous requirement.
			\item \textbf{Why A is wrong:} Use test results are not technical spreadsheets
			\item \textbf{Why C is wrong:} Results go to senior management
			\item \textbf{Why D is wrong:} Results are documented in some form
		\end{itemize}
		
		\textcolor{myred}{\textbf{EXAM TRAP:}} Use test results go to \textcolor{myblue}{senior management} - not technical reports!
	\end{frame}
	
	\begin{frame}{Q47: Model Use Limits - Tutorial}
		\fontsize{6.5}{7.5}\selectfont
		\textbf{QUESTION: Why should there be limits on model use?}
		
		\vspace{0.1cm}
		\textbf{OPTIONS:}
		\begin{enumerate}
			\item To make models harder to use
			\item Because every QRM is based on underlying assumptions with limitations that should be clearly stated and communicated
			\item To reduce the number of models
			\item To increase model risk
		\end{enumerate}
		
		\vspace{0.1cm}
		\textcolor{myblue}{\textbf{ANSWER: B}}
		
		\vspace{0.1cm}
		\textbf{DETAILED EXPLANATION:}
		\begin{itemize}
			\item \textbf{Why B is correct:} Every quantitative risk model is based on underlying assumptions that should be clearly stated in the documentation and communicated to model users. However, it is important to identify model limitations and weaknesses and to establish limits on the usage of models to prevent misuse.
			\item \textbf{Why A is wrong:} The purpose is not to make models harder to use
			\item \textbf{Why C is wrong:} Limits are about appropriate use, not reducing count
			\item \textbf{Why D is wrong:} Limits reduce risk, not increase it
		\end{itemize}
		
		\textcolor{myred}{\textbf{EXAM TRAP:}} Model limits prevent \textcolor{myblue}{misuse beyond intended application}!
	\end{frame}
	
	\begin{frame}{Q48: Three Guidelines - Tutorial}
		\fontsize{6.5}{7.5}\selectfont
		\textbf{QUESTION: What are the three general guidelines for risk modeling culture?}
		
		\vspace{0.1cm}
		\textbf{OPTIONS:}
		\begin{enumerate}
			\item Speed, cost, and efficiency
			\item Awareness, transparency, and experience
			\item Development, validation, and deployment
			\item Data, models, and reports
		\end{enumerate}
		
		\vspace{0.1cm}
		\textcolor{myblue}{\textbf{ANSWER: B}}
		
		\vspace{0.1cm}
		\textbf{DETAILED EXPLANATION:}
		\begin{itemize}
			\item \textbf{Why B is correct:} Three general guidelines can support management, modelers, and validators: Awareness (be aware of limitations and assumptions), Transparency (transparently communicate assumptions, limitations, and documentation), and Experience (learn from past modeling endeavors and apply relevant lessons).
			\item \textbf{Why A is wrong:} These are operational metrics, not governance guidelines
			\item \textbf{Why C is wrong:} These are lifecycle phases
			\item \textbf{Why D is wrong:} These are technical components
		\end{itemize}
		
		\textcolor{myred}{\textbf{EXAM TRAP:}} Risk modeling culture guidelines are \textcolor{myblue}{awareness, transparency, and experience}!
	\end{frame}
	
	\begin{frame}{Q49: Model Change Process - Tutorial}
		\fontsize{6.5}{7.5}\selectfont
		\textbf{QUESTION: What is the first step in the model change management process?}
		
		\vspace{0.1cm}
		\textbf{OPTIONS:}
		\begin{enumerate}
			\item Implementation
			\item Change requests
			\item Validation
			\item Documentation
		\end{enumerate}
		
		\vspace{0.1cm}
		\textcolor{myblue}{\textbf{ANSWER: B}}
		
		\vspace{0.1cm}
		\textbf{DETAILED EXPLANATION:}
		\begin{itemize}
			\item \textbf{Why B is correct:} The change-management process includes activities such as change requests, impact assessment, documentation, validation, and approval. Change requests initiate the process.
			\item \textbf{Why A is wrong:} Implementation comes after approval
			\item \textbf{Why C is wrong:} Validation comes after impact assessment
			\item \textbf{Why D is wrong:} Documentation occurs throughout the process
		\end{itemize}
		
		\textcolor{myred}{\textbf{EXAM TRAP:}} Model changes start with \textcolor{myblue}{change requests} and follow structured process!
	\end{frame}
	
	\begin{frame}{Q50: AI/ML Review Differences - Tutorial}
		\fontsize{6.5}{7.5}\selectfont
		\textbf{QUESTION: How does the review framework for AI/ML models differ from traditional models?}
		
		\vspace{0.1cm}
		\textbf{OPTIONS:}
		\begin{enumerate}
			\item AI/ML models never need review
			\item AI/ML models require more frequent monitoring, broader stakeholder involvement, and focus on explainability
			\item AI/ML models use the same review framework as traditional models
			\item AI/ML models only need review once
		\end{enumerate}
		
		\vspace{0.1cm}
		\textcolor{myblue}{\textbf{ANSWER: B}}
		
		\vspace{0.1cm}
		\textbf{DETAILED EXPLANATION:}
		\begin{itemize}
			\item \textbf{Why B is correct:} The review framework differs in three notable ways: (1) Frequency - more frequent monitoring needed due to complexity and data drift, (2) Validation stakeholders - wider range of stakeholders affected, and (3) Validation content - focus on explainability, benchmarking, and bias detection.
			\item \textbf{Why A is wrong:} AI/ML models require rigorous review
			\item \textbf{Why C is wrong:} The framework must adapt to AI/ML characteristics
			\item \textbf{Why D is wrong:} Continuous review is needed
		\end{itemize}
		
		\textcolor{myred}{\textbf{EXAM TRAP:}} AI/ML review requires \textcolor{myblue}{more frequent monitoring, broader stakeholders, and explainability focus}!
	\end{frame}
	
	% ============================================================
	% SECTION 5: MODEL VALIDATION - QUESTIONS 51-65
	% ============================================================
	
	\begin{frame}{Q51: Model Validation Definition - Tutorial}
		\fontsize{6.5}{7.5}\selectfont
		\textbf{QUESTION: What is model validation?}
		
		\vspace{0.1cm}
		\textbf{OPTIONS:}
		\begin{enumerate}
			\item The process of developing a model
			\item The process of rigorously testing the model, inputs, outputs, governance, etc. to try to identify the extent of residual model risk and assess mitigating controls
			\item The process of deploying a model
			\item The process of documenting a model
		\end{enumerate}
		
		\vspace{0.1cm}
		\textcolor{myblue}{\textbf{ANSWER: B}}
		
		\vspace{0.1cm}
		\textbf{DETAILED EXPLANATION:}
		\begin{itemize}
			\item \textbf{Why B is correct:} Model validation is the process of rigorously testing the model, inputs, outputs, governance, etc. to try to identify the extent of residual model risk and assess mitigating controls.
			\item \textbf{Why A is wrong:} Development and validation are separate processes
			\item \textbf{Why C is wrong:} Deployment comes after validation
			\item \textbf{Why D is wrong:} Documentation is part of development
		\end{itemize}
		
		\textcolor{myred}{\textbf{EXAM TRAP:}} Model validation is \textcolor{myblue}{rigorous testing to identify residual risk}!
	\end{frame}
	
	\begin{frame}{Q52: Initial Validation - Tutorial}
		\fontsize{6.5}{7.5}\selectfont
		\textbf{QUESTION: When is initial validation performed?}
		
		\vspace{0.1cm}
		\textbf{OPTIONS:}
		\begin{enumerate}
			\item Only after the model has been in production for a year
			\item Prior to implementation and usage
			\item Only when the model fails
			\item Never - models don't need initial validation
		\end{enumerate}
		
		\vspace{0.1cm}
		\textcolor{myblue}{\textbf{ANSWER: B}}
		
		\vspace{0.1cm}
		\textbf{DETAILED EXPLANATION:}
		\begin{itemize}
			\item \textbf{Why B is correct:} Once a model is identified, it is inventoried and scheduled for an initial baseline validation, which occurs prior to implementation and usage.
			\item \textbf{Why A is wrong:} Initial validation occurs before production
			\item \textbf{Why C is wrong:} Validation is proactive, not reactive
			\item \textbf{Why D is wrong:} Initial validation is essential
		\end{itemize}
		
		\textcolor{myred}{\textbf{EXAM TRAP:}} Initial validation occurs \textcolor{myblue}{before implementation and usage}!
	\end{frame}
	
	\begin{frame}{Q53: Periodic Validation - Tutorial}
		\fontsize{6.5}{7.5}\selectfont
		\textbf{QUESTION: What is the primary objective of ongoing (periodic) validation?}
		
		\vspace{0.1cm}
		\textbf{OPTIONS:}
		\begin{enumerate}
			\item To start over with a new model
			\item To observe whether the model remains aligned with its intended purpose, assumptions remain valid, data is still appropriate, and performance monitoring indicates expected performance
			\item To replace the model entirely
			\item To ignore model performance
		\end{enumerate}
		
		\vspace{0.1cm}
		\textcolor{myblue}{\textbf{ANSWER: B}}
		
		\vspace{0.1cm}
		\textbf{DETAILED EXPLANATION:}
		\begin{itemize}
			\item \textbf{Why B is correct:} The primary objective of ongoing or periodic validation is to observe whether the model remains aligned with its intended purpose, the assumptions remain valid, the data are still appropriate, and performance monitoring indicates that the model continues to perform as expected.
			\item \textbf{Why A is wrong:} Ongoing validation is about monitoring, not restarting
			\item \textbf{Why C is wrong:} Replacement is only one option
			\item \textbf{Why D is wrong:} Performance must be monitored
		\end{itemize}
		
		\textcolor{myred}{\textbf{EXAM TRAP:}} Ongoing validation ensures \textcolor{myblue}{model remains fit for purpose over time}!
	\end{frame}
	
	\begin{frame}{Q54: Validation Frequency - Tutorial}
		\fontsize{6.5}{7.5}\selectfont
		\textbf{QUESTION: How should the frequency and intensity of validation be determined?}
		
		\vspace{0.1cm}
		\textbf{OPTIONS:}
		\begin{enumerate}
			\item All models should be validated at the same frequency
			\item Based on the risk ranking of the model
			\item Only when regulators request
			\item Only when the model fails
		\end{enumerate}
		
		\vspace{0.1cm}
		\textcolor{myblue}{\textbf{ANSWER: B}}
		
		\vspace{0.1cm}
		\textbf{DETAILED EXPLANATION:}
		\begin{itemize}
			\item \textbf{Why B is correct:} The frequency and intensity of validation activities should be determined based on the risk ranking of the model. High-risk models may have a periodic baseline revalidation every year or more frequently, whereas lower-risk models may be revalidated on a two- or three-year frequency.
			\item \textbf{Why A is wrong:} Different risk tiers require different frequencies
			\item \textbf{Why C is wrong:} Validation should be proactive, not just reactive
			\item \textbf{Why D is wrong:} Don't wait for failure to validate
		\end{itemize}
		
		\textcolor{myred}{\textbf{EXAM TRAP:}} Validation frequency is \textcolor{myblue}{based on model risk ranking} - higher risk = more frequent!
	\end{frame}
	
	\begin{frame}{Q55: Validation Activities - Tutorial}
		\fontsize{6.5}{7.5}\selectfont
		\textbf{QUESTION: Which of the following is NOT a model validation activity?}
		
		\vspace{0.1cm}
		\textbf{OPTIONS:}
		\begin{enumerate}
			\item Assessment of data sources
			\item Review of conceptual soundness
			\item Model development
			\item Review of quality of documentation
		\end{enumerate}
		
		\vspace{0.1cm}
		\textcolor{myblue}{\textbf{ANSWER: C}}
		
		\vspace{0.1cm}
		\textbf{DETAILED EXPLANATION:}
		\begin{itemize}
			\item \textbf{Why C is correct:} Model development is a separate activity that precedes validation. Validation activities include assessment of data sources, assumptions, limitations, implementation, conceptual soundness, performance monitoring, documentation, controls, outcomes analysis, and testing.
			\item \textbf{Why A is wrong:} Data assessment IS a validation activity
			\item \textbf{Why B is wrong:} Conceptual soundness review IS a validation activity
			\item \textbf{Why D is wrong:} Documentation review IS a validation activity
		\end{itemize}
		
		\textcolor{myred}{\textbf{EXAM TRAP:}} Validation is \textcolor{myblue}{separate from development} - not the same activity!
	\end{frame}
	
	\begin{frame}{Q56: Model Validation Terminology - Tutorial}
		\fontsize{6.5}{7.5}\selectfont
		\textbf{QUESTION: Which is NOT an alternative term for model validation?}
		
		\vspace{0.1cm}
		\textbf{OPTIONS:}
		\begin{enumerate}
			\item Model review
			\item Model QA (quality assurance)
			\item Model development
			\item Model vetting
		\end{enumerate}
		
		\vspace{0.1cm}
		\textcolor{myblue}{\textbf{ANSWER: C}}
		
		\vspace{0.1cm}
		\textbf{DETAILED EXPLANATION:}
		\begin{itemize}
			\item \textbf{Why C is correct:} Model development is the process of creating the model, not a term for validation. Alternative terms include model review, model QA (quality assurance), model evaluation, and model vetting.
			\item \textbf{Why A is wrong:} Model review IS an alternative term
			\item \textbf{Why B is wrong:} Model QA IS an alternative term
			\item \textbf{Why D is wrong:} Model vetting IS an alternative term
		\end{itemize}
		
		\textcolor{myred}{\textbf{EXAM TRAP:}} Validation is NOT development - different terms for different activities!
	\end{frame}
	
	\begin{frame}{Q57: Recursive Validation - Tutorial}
		\fontsize{6.5}{7.5}\selectfont
		\textbf{QUESTION: What does "recursive validation" mean?}
		
		\vspace{0.1cm}
		\textbf{OPTIONS:}
		\begin{enumerate}
			\item Validation that never ends
			\item Validation activities should not be immune to critical examination
			\item Validation that only happens once
			\item Validation that ignores all feedback
		\end{enumerate}
		
		\vspace{0.1cm}
		\textcolor{myblue}{\textbf{ANSWER: B}}
		
		\vspace{0.1cm}
		\textbf{DETAILED EXPLANATION:}
		\begin{itemize}
			\item \textbf{Why B is correct:} Validation activities should be "recursive," meaning that they should not be immune to critical examination. It is essential to ensure that the validation process does not rely on defective or inappropriate data and remains within the model's intended application domain.
			\item \textbf{Why A is wrong:} Recursive means open to challenge, not never-ending
			\item \textbf{Why C is wrong:} Validation is ongoing, not one-time
			\item \textbf{Why D is wrong:} Feedback is essential
		\end{itemize}
		
		\textcolor{myred}{\textbf{EXAM TRAP:}} Recursive validation means \textcolor{myblue}{validation itself must be subject to challenge}!
	\end{frame}
	
	\begin{frame}{Q58: No "Valid" Model - Tutorial}
		\fontsize{6.5}{7.5}\selectfont
		\textbf{QUESTION: Why is the goal not to test for a "valid" model?}
		
		\vspace{0.1cm}
		\textbf{OPTIONS:}
		\begin{enumerate}
			\item Because all models are perfectly valid
			\item Because models are simplified representations of reality under certain assumed conditions - there is no perfect model
			\item Because validation is unnecessary
			\item Because models never need testing
		\end{enumerate}
		
		\vspace{0.1cm}
		\textcolor{myblue}{\textbf{ANSWER: B}}
		
		\vspace{0.1cm}
		\textbf{DETAILED EXPLANATION:}
		\begin{itemize}
			\item \textbf{Why B is correct:} Because models are simplified representations of reality under certain assumed conditions, there is no perfect model. The goal of the validation exercise is thus not to test for a "valid" or fully validated model, but rather to subject the model to a series of attempts to invalidate it.
			\item \textbf{Why A is wrong:} No model is perfectly valid
			\item \textbf{Why C is wrong:} Validation is essential
			\item \textbf{Why D is wrong:} Models need rigorous testing
		\end{itemize}
		
		\textcolor{myred}{\textbf{EXAM TRAP:}} Goal is to \textcolor{myblue}{attempt to invalidate} - not prove validity!
	\end{frame}
	
	\begin{frame}{Q59: Effective Challenge - Tutorial}
		\fontsize{6.5}{7.5}\selectfont
		\textbf{QUESTION: According to SR 11-7, what does "effective challenge" require?}
		
		\vspace{0.1cm}
		\textbf{OPTIONS:}
		\begin{enumerate}
			\item Only technical expertise
			\item Critical analysis by objective and informed parties who can identify assumptions, limitations, and weaknesses
			\item Only regulatory approval
			\item Only management approval
		\end{enumerate}
		
		\vspace{0.1cm}
		\textcolor{myblue}{\textbf{ANSWER: B}}
		
		\vspace{0.1cm}
		\textbf{DETAILED EXPLANATION:}
		\begin{itemize}
			\item \textbf{Why B is correct:} The Fed's SR 11-7 speaks of the "effective challenge" of models. This should encompass a critical analysis by objective and informed parties that are able to identify crucial model assumptions and limitations, and to spot and communicate relevant model weaknesses.
			\item \textbf{Why A is wrong:} Technical expertise alone is insufficient
			\item \textbf{Why C is wrong:} Regulatory approval is external
			\item \textbf{Why D is wrong:} Management approval is internal
		\end{itemize}
		
		\textcolor{myred}{\textbf{EXAM TRAP:}} Effective challenge requires \textcolor{myblue}{objective, informed, critical analysis}!
	\end{frame}
	
	\begin{frame}{Q60: GenAI Validation - Special Considerations - Tutorial}
		\fontsize{6.5}{7.5}\selectfont
		\textbf{QUESTION: Which is NOT a special consideration for GenAI validation?}
		
		\vspace{0.1cm}
		\textbf{OPTIONS:}
		\begin{enumerate}
			\item Open-ended output
			\item Hallucination
			\item Simple and straightforward validation like traditional models
			\item Multi-modal issues
		\end{enumerate}
		
		\vspace{0.1cm}
		\textcolor{myblue}{\textbf{ANSWER: C}}
		
		\vspace{0.1cm}
		\textbf{DETAILED EXPLANATION:}
		\begin{itemize}
			\item \textbf{Why C is correct:} GenAI validation is NOT simple and straightforward. It requires special considerations including open-ended output, different performance metrics, data drift, hallucinations, discrimination and bias, multi-modal issues, interpretability challenges, rapid technological change, third-party models, and significant computational resources.
			\item \textbf{Why A is wrong:} Open-ended output IS a special consideration
			\item \textbf{Why B is wrong:} Hallucination IS a special consideration
			\item \textbf{Why D is wrong:} Multi-modal issues ARE special considerations
		\end{itemize}
		
		\textcolor{myred}{\textbf{EXAM TRAP:}} GenAI validation is \textcolor{myblue}{complex with unique challenges} - not straightforward!
	\end{frame}
	
	\begin{frame}{Q61: Hallucination in GenAI - Tutorial}
		\fontsize{6.5}{7.5}\selectfont
		\textbf{QUESTION: What is a key concern when validating GenAI models regarding hallucinations?}
		
		\vspace{0.1cm}
		\textbf{OPTIONS:}
		\begin{enumerate}
			\item GenAI models never hallucinate
			\item GenAI models have a tendency to hallucinate or give responses that are factually incorrect but apparently very convincing
			\item Hallucinations are always beneficial
			\item Hallucinations are easy to detect
		\end{enumerate}
		
		\vspace{0.1cm}
		\textcolor{myblue}{\textbf{ANSWER: B}}
		
		\vspace{0.1cm}
		\textbf{DETAILED EXPLANATION:}
		\begin{itemize}
			\item \textbf{Why B is correct:} GenAI models have a tendency to hallucinate or give responses that are factually incorrect but apparently very convincing. This makes validation particularly challenging as outputs may appear credible but be false.
			\item \textbf{Why A is wrong:} Hallucination is a well-known GenAI issue
			\item \textbf{Why C is wrong:} Hallucinations are harmful, not beneficial
			\item \textbf{Why D is wrong:} Hallucinations are often difficult to detect
		\end{itemize}
		
		\textcolor{myred}{\textbf{EXAM TRAP:}} Hallucinations in GenAI are \textcolor{myblue}{factually incorrect but appear convincing}!
	\end{frame}
	
	\begin{frame}{Q62: Model Design Motivation - Tutorial}
		\fontsize{6.5}{7.5}\selectfont
		\textbf{QUESTION: Which is NOT a motivation for developing a new QRM?}
		
		\vspace{0.1cm}
		\textbf{OPTIONS:}
		\begin{enumerate}
			\item Changing conditions (new markets, new products)
			\item Regulatory pressure
			\item Recent crisis
			\item To make models more complex without reason
		\end{enumerate}
		
		\vspace{0.1cm}
		\textcolor{myblue}{\textbf{ANSWER: D}}
		
		\vspace{0.1cm}
		\textbf{DETAILED EXPLANATION:}
		\begin{itemize}
			\item \textbf{Why D is correct:} Making models more complex without reason is not a valid motivation. Valid motivations include changing conditions, regulatory pressure, recent crises, internal concerns, and innovation/business growth.
			\item \textbf{Why A is wrong:} Changing conditions IS a valid motivation
			\item \textbf{Why B is wrong:} Regulatory pressure IS a valid motivation
			\item \textbf{Why C is wrong:} Recent crisis IS a valid motivation
		\end{itemize}
		
		\textcolor{myred}{\textbf{EXAM TRAP:}} Model development should be driven by \textcolor{myblue}{business need, not complexity for its own sake}!
	\end{frame}
	
	\begin{frame}{Q63: Discretization - Tutorial}
		\fontsize{6.5}{7.5}\selectfont
		\textbf{QUESTION: What is discretization in modeling?}
		
		\vspace{0.1cm}
		\textbf{OPTIONS:}
		\begin{enumerate}
			\item Using only discrete values without considering continuous ones
			\item Dividing time intervals into days, months, or years, and rounding monetary gains or losses to discrete measures
			\item Using only continuous variables
			\item Ignoring time intervals
		\end{enumerate}
		
		\vspace{0.1cm}
		\textcolor{myblue}{\textbf{ANSWER: B}}
		
		\vspace{0.1cm}
		\textbf{DETAILED EXPLANATION:}
		\begin{itemize}
			\item \textbf{Why B is correct:} In discretization, time intervals are divided into days, months, or years, and monetary gains or losses are rounded to cents or other discrete measures. The discretization should accurately reflect the underlying physical or mathematical principles.
			\item \textbf{Why A is wrong:} Discretization involves dividing continuous values
			\item \textbf{Why C is wrong:} Discretization is the opposite of using only continuous
			\item \textbf{Why D is wrong:} Time intervals are considered
		\end{itemize}
		
		\textcolor{myred}{\textbf{EXAM TRAP:}} Discretization involves \textcolor{myblue}{dividing continuous variables into discrete measures}!
	\end{frame}
	
	\begin{frame}{Q64: Approximations - Tutorial}
		\fontsize{6.5}{7.5}\selectfont
		\textbf{QUESTION: What is a key characteristic of approximations in modeling?}
		
		\vspace{0.1cm}
		\textbf{OPTIONS:}
		\begin{enumerate}
			\item They can be improved indefinitely
			\item Choosing an approximation fixes the precision, and there is no room for improvement except by selecting a different approximation
			\item They always produce exact results
			\item They are never used in risk modeling
		\end{enumerate}
		
		\vspace{0.1cm}
		\textcolor{myblue}{\textbf{ANSWER: B}}
		
		\vspace{0.1cm}
		\textbf{DETAILED EXPLANATION:}
		\begin{itemize}
			\item \textbf{Why B is correct:} Approximation involves replacing complex or hard-to-evaluate model components with alternative methods that produce similar results more conveniently. However, choosing an approximation fixes the precision, and there is no room for improvement except by selecting a different approximation.
			\item \textbf{Why A is wrong:} Precision is fixed by the approximation choice
			\item \textbf{Why C is wrong:} Approximations are not exact
			\item \textbf{Why D is wrong:} Approximations are commonly used
		\end{itemize}
		
		\textcolor{myred}{\textbf{EXAM TRAP:}} Approximations \textcolor{myblue}{fix precision} - can only improve by choosing different approximation!
	\end{frame}
	
	\begin{frame}{Q65: Numerical Evaluation - Tutorial}
		\fontsize{6.5}{7.5}\selectfont
		\textbf{QUESTION: How does numerical evaluation differ from approximation?}
		
		\vspace{0.1cm}
		\textbf{OPTIONS:}
		\begin{enumerate}
			\item Numerical evaluation is less precise
			\item Numerical evaluation allows for potential arbitrary precision through parameters controlling order, iterations, step size, and grid size
			\item Numerical evaluation is always simpler
			\item Numerical evaluation never requires testing
		\end{enumerate}
		
		\vspace{0.1cm}
		\textcolor{myblue}{\textbf{ANSWER: B}}
		
		\vspace{0.1cm}
		\textbf{DETAILED EXPLANATION:}
		\begin{itemize}
			\item \textbf{Why B is correct:} In contrast to approximation, numerical evaluation means evaluation of a model component with the potential for arbitrary precision. Therefore, methods for numerical evaluation are frequently equipped with parameters allowing control of precision (order, number of iterations, step size, grid size).
			\item \textbf{Why A is wrong:} Numerical evaluation can be more precise
			\item \textbf{Why C is wrong:} It can be more complex
			\item \textbf{Why D is wrong:} Testing is essential
		\end{itemize}
		
		\textcolor{myred}{\textbf{EXAM TRAP:}} Numerical evaluation allows \textcolor{myblue}{arbitrary precision through parameter control}!
	\end{frame}
	
	% ============================================================
	% SECTION 6: IMPLEMENTATION AND MISINTERPRETATION - QUESTIONS 66-75
	% ============================================================
	
	\begin{frame}{Q66: Implementation Tasks - Tutorial}
		\fontsize{6.5}{7.5}\selectfont
		\textbf{QUESTION: What are the three model implementation tasks?}
		
		\vspace{0.1cm}
		\textbf{OPTIONS:}
		\begin{enumerate}
			\item Design, testing, deployment
			\item Model design, core implementation, system implementation
			\item Training, validation, testing
			\item Development, validation, monitoring
		\end{enumerate}
		
		\vspace{0.1cm}
		\textcolor{myblue}{\textbf{ANSWER: B}}
		
		\vspace{0.1cm}
		\textbf{DETAILED EXPLANATION:}
		\begin{itemize}
			\item \textbf{Why B is correct:} The three implementation tasks are: (1) Model design - designed and documented to facilitate translation into computer code, (2) Core implementation - inner workings implemented as a black box with specified interfaces, and (3) System implementation - model core integrated into a system that provides user interfaces, collects input data, schedules computations, and processes output data.
			\item \textbf{Why A is wrong:} Testing and deployment are separate phases
			\item \textbf{Why C is wrong:} These are data splits
			\item \textbf{Why D is wrong:} These are lifecycle phases
		\end{itemize}
		
		\textcolor{myred}{\textbf{EXAM TRAP:}} Implementation tasks are \textcolor{myblue}{model design, core implementation, and system implementation}!
	\end{frame}
	
	\begin{frame}{Q67: Core vs System Implementation - Tutorial}
		\fontsize{6.5}{7.5}\selectfont
		\textbf{QUESTION: What is the difference between core implementation and system implementation?}
		
		\vspace{0.1cm}
		\textbf{OPTIONS:}
		\begin{enumerate}
			\item There is no difference
			\item Core implementation is the inner workings; system implementation integrates the core into a system with user interfaces, data collection, scheduling, and output processing
			\item Core implementation is easier than system implementation
			\item System implementation is done first
		\end{enumerate}
		
		\vspace{0.1cm}
		\textcolor{myblue}{\textbf{ANSWER: B}}
		
		\vspace{0.1cm}
		\textbf{DETAILED EXPLANATION:}
		\begin{itemize}
			\item \textbf{Why B is correct:} Core implementation involves the inner workings of the model as a black box with specified interfaces. System implementation integrates the model core into a system that provides user interfaces, collects input data and parameters, schedules computations, feeds the core, processes output data, and keeps a history of computations.
			\item \textbf{Why A is wrong:} They are distinct tasks with different requirements
			\item \textbf{Why C is wrong:} Both have different challenges
			\item \textbf{Why D is wrong:} Core implementation typically comes first
		\end{itemize}
		
		\textcolor{myred}{\textbf{EXAM TRAP:}} Core = \textcolor{myblue}{inner workings}; System = \textcolor{myblue}{integration with user interface and data flows}!
	\end{frame}
	
	\begin{frame}{Q68: Model Adaptation - Tutorial}
		\fontsize{6.5}{7.5}\selectfont
		\textbf{QUESTION: What are the three model adaptation tasks?}
		
		\vspace{0.1cm}
		\textbf{OPTIONS:}
		\begin{enumerate}
			\item Model adaptation, core adaptation, system adaptation
			\item Only model adaptation
			\item Only core adaptation
			\item Only system adaptation
		\end{enumerate}
		
		\vspace{0.1cm}
		\textcolor{myblue}{\textbf{ANSWER: A}}
		
		\vspace{0.1cm}
		\textbf{DETAILED EXPLANATION:}
		\begin{itemize}
			\item \textbf{Why A is correct:} Because no model lasts forever, periodic reviews may detect weaknesses that demand adaptation: (1) Model adaptation - model and documentation adapted, (2) Core adaptation - model core and interfaces adapted, and (3) System adaptation - system around the model core adapted or replaced.
			\item \textbf{Why B is wrong:} There are three types of adaptation
			\item \textbf{Why C is wrong:} There are three types of adaptation
			\item \textbf{Why D is wrong:} There are three types of adaptation
		\end{itemize}
		
		\textcolor{myred}{\textbf{EXAM TRAP:}} Adaptation involves \textcolor{myblue}{model, core, and system} adaptation!
	\end{frame}
	
	\begin{frame}{Q69: Running a Model - Tutorial}
		\fontsize{6.5}{7.5}\selectfont
		\textbf{QUESTION: What are the two ways a model can be run?}
		
		\vspace{0.1cm}
		\textbf{OPTIONS:}
		\begin{enumerate}
			\item Fast and slow
			\item Scheduled runs and ad hoc runs
			\item Manual and automatic
			\item Internal and external
		\end{enumerate}
		
		\vspace{0.1cm}
		\textcolor{myblue}{\textbf{ANSWER: B}}
		
		\vspace{0.1cm}
		\textbf{DETAILED EXPLANATION:}
		\begin{itemize}
			\item \textbf{Why B is correct:} A desirable aspect of a model is its ability to be run in two ways: scheduled runs (regular numbers with minimal manual intervention) and ad hoc runs (additional information on short notice, usually requiring manual intervention).
			\item \textbf{Why A is wrong:} Speed is not the classification
			\item \textbf{Why C is wrong:} Both scheduled and ad hoc can be automatic or manual
			\item \textbf{Why D is wrong:} Location is not the classification
		\end{itemize}
		
		\textcolor{myred}{\textbf{EXAM TRAP:}} Model runs can be \textcolor{myblue}{scheduled (regular) or ad hoc (on-demand)}!
	\end{frame}
	
	\begin{frame}{Q70: Misinterpretation Sources - Tutorial}
		\fontsize{6.5}{7.5}\selectfont
		\textbf{QUESTION: Which is NOT a source of model result misinterpretation?}
		
		\vspace{0.1cm}
		\textbf{OPTIONS:}
		\begin{enumerate}
			\item Lack of comprehension
			\item Neglecting model limitations
			\item Perfect understanding of all assumptions
			\item Overemphasis on point estimates
		\end{enumerate}
		
		\vspace{0.1cm}
		\textcolor{myblue}{\textbf{ANSWER: C}}
		
		\vspace{0.1cm}
		\textbf{DETAILED EXPLANATION:}
		\begin{itemize}
			\item \textbf{Why C is correct:} Perfect understanding of all assumptions is not a source of misinterpretation - it's the ideal state. Sources include lack of comprehension, neglecting limitations, overemphasis on point estimates, failure to consider qualitative factors, ignoring decision context, and confirmation bias.
			\item \textbf{Why A is wrong:} Lack of comprehension IS a source
			\item \textbf{Why B is wrong:} Neglecting limitations IS a source
			\item \textbf{Why D is wrong:} Overemphasis on point estimates IS a source
		\end{itemize}
		
		\textcolor{myred}{\textbf{EXAM TRAP:}} Misinterpretation arises from \textcolor{myblue}{lack of understanding, not perfect understanding}!
	\end{frame}
	
	\begin{frame}{Q71: Confirmation Bias - Tutorial}
		\fontsize{6.5}{7.5}\selectfont
		\textbf{QUESTION: What is confirmation bias in interpreting model results?}
		
		\vspace{0.1cm}
		\textbf{OPTIONS:}
		\begin{enumerate}
			\item Always interpreting results correctly
			\item Users interpreting results to align with their preconceived beliefs or desired outcomes
			\item Ignoring all results
			\item Always questioning results
		\end{enumerate}
		
		\vspace{0.1cm}
		\textcolor{myblue}{\textbf{ANSWER: B}}
		
		\vspace{0.1cm}
		\textbf{DETAILED EXPLANATION:}
		\begin{itemize}
			\item \textbf{Why B is correct:} Confirmation bias occurs when users may unintentionally interpret model results in a way that aligns with their preconceived beliefs or desired outcomes. Confirmation bias can cloud judgment and compromise the objective interpretation of the model results.
			\item \textbf{Why A is wrong:} This is the opposite of confirmation bias
			\item \textbf{Why C is wrong:} Ignoring results is not confirmation bias
			\item \textbf{Why D is wrong:} Questioning results is good practice
		\end{itemize}
		
		\textcolor{myred}{\textbf{EXAM TRAP:}} Confirmation bias means \textcolor{myblue}{interpreting results to fit preconceived beliefs}!
	\end{frame}
	
	\begin{frame}{Q72: Preventing Misinterpretation - Tutorial}
		\fontsize{6.5}{7.5}\selectfont
		\textbf{QUESTION: How can organizations foster proper interpretation of model results?}
		
		\vspace{0.1cm}
		\textbf{OPTIONS:}
		\begin{enumerate}
			\item By ignoring all results
			\item Through training, clear communication, open dialogue, and a supportive culture
			\item By only using automated interpretation
			\item By never questioning results
		\end{enumerate}
		
		\vspace{0.1cm}
		\textcolor{myblue}{\textbf{ANSWER: B}}
		
		\vspace{0.1cm}
		\textbf{DETAILED EXPLANATION:}
		\begin{itemize}
			\item \textbf{Why B is correct:} Fostering a culture that encourages open dialogue, critical thinking, and cross-functional collaboration can enhance the interpretation process and guard against misinterpretation biases. This requires training, clear communication, and a supportive culture.
			\item \textbf{Why A is wrong:} Ignoring results is not a solution
			\item \textbf{Why C is wrong:} Human interpretation remains essential
			\item \textbf{Why D is wrong:} Results should be questioned and critically examined
		\end{itemize}
		
		\textcolor{myred}{\textbf{EXAM TRAP:}} Proper interpretation requires \textcolor{myblue}{training, communication, and supportive culture}!
	\end{frame}
	
	\begin{frame}{Q73: AI/ML Governance Recommendations - Tutorial}
		\fontsize{6.5}{7.5}\selectfont
		\textbf{QUESTION: What is the first recommendation for establishing AI/ML governance frameworks?}
		
		\vspace{0.1cm}
		\textbf{OPTIONS:}
		\begin{enumerate}
			\item Start from scratch
			\item Begin with existing model governance frameworks
			\item Ignore all existing frameworks
			\item Only use vendor frameworks
		\end{enumerate}
		
		\vspace{0.1cm}
		\textcolor{myblue}{\textbf{ANSWER: B}}
		
		\vspace{0.1cm}
		\textbf{DETAILED EXPLANATION:}
		\begin{itemize}
			\item \textbf{Why B is correct:} Begin with existing model governance frameworks. Even though AI/ML introduces new challenges for model risk management, enhanced model governance frameworks should not start from scratch. Financial institutions and regulators have gained much experience in model validation in recent years that could serve as strong basis.
			\item \textbf{Why A is wrong:} Don't start from scratch - build on existing frameworks
			\item \textbf{Why C is wrong:} Existing frameworks provide valuable foundation
			\item \textbf{Why D is wrong:} Vendor frameworks should complement, not replace, internal frameworks
		\end{itemize}
		
		\textcolor{myred}{\textbf{EXAM TRAP:}} AI/ML governance should \textcolor{myblue}{build on existing model governance frameworks}!
	\end{frame}
	
	\begin{frame}{Q74: New Role of Data in AI/ML - Tutorial}
		\fontsize{6.5}{7.5}\selectfont
		\textbf{QUESTION: How does AI/ML change the role of data in model governance?}
		
		\vspace{0.1cm}
		\textbf{OPTIONS:}
		\begin{enumerate}
			\item Data becomes less important
			\item Data's role becomes more important - AI/ML is "model free" and everything depends only on the data
			\item Data becomes irrelevant
			\item Data is only needed for training
		\end{enumerate}
		
		\vspace{0.1cm}
		\textcolor{myblue}{\textbf{ANSWER: B}}
		
		\vspace{0.1cm}
		\textbf{DETAILED EXPLANATION:}
		\begin{itemize}
			\item \textbf{Why B is correct:} The new paradigm suggests that AI/ML is "model free," and everything depends only on the data. Though that may not be literally true, the more important role of data needs to be addressed within model governance frameworks. This addresses issues related to bias, overfitting, drift, and harmful output.
			\item \textbf{Why A is wrong:} Data becomes MORE important, not less
			\item \textbf{Why C is wrong:} Data is central to AI/ML
			\item \textbf{Why D is wrong:} Data's role extends beyond just training
		\end{itemize}
		
		\textcolor{myred}{\textbf{EXAM TRAP:}} In AI/ML, \textcolor{myblue}{data plays an even more central role} than in traditional models!
	\end{frame}
	
	\begin{frame}{Q75: BIS Survey - Data Governance - Tutorial}
		\fontsize{6.5}{7.5}\selectfont
		\textbf{QUESTION: According to the BIS survey, what is a common solution for data governance?}
		
		\vspace{0.1cm}
		\textbf{OPTIONS:}
		\begin{enumerate}
			\item No data governance is needed
			\item Establishing comprehensive data management systems, documenting data through metadata, and maintaining asset inventories
			\item Only using cloud storage
			\item Only using spreadsheets
		\end{enumerate}
		
		\vspace{0.1cm}
		\textcolor{myblue}{\textbf{ANSWER: B}}
		
		\vspace{0.1cm}
		\textbf{DETAILED EXPLANATION:}
		\begin{itemize}
			\item \textbf{Why B is correct:} Common solutions typically include establishing comprehensive data management systems, documenting data (e.g., through metadata on the sources, version and curation), maintaining asset inventories (e.g., data catalogues) and metadata registries, using standards and enabling efficient data exploration for both humans and machines.
			\item \textbf{Why A is wrong:} Data governance is essential
			\item \textbf{Why C is wrong:} Cloud storage is just one tool
			\item \textbf{Why D is wrong:} Spreadsheets are insufficient
		\end{itemize}
		
		\textcolor{myred}{\textbf{EXAM TRAP:}} BIS survey shows data governance requires \textcolor{myblue}{comprehensive systems, metadata, and inventories}!
	\end{frame}
	
	% ============================================================
	% SECTION 7: GENAI GOVERNANCE - QUESTIONS 76-85
	% ============================================================
	
	\begin{frame}{Q76: GenAI Governance Challenges - Tutorial}
		\fontsize{6.5}{7.5}\selectfont
		\textbf{QUESTION: Why is GenAI governance harder than traditional model governance?}
		
		\vspace{0.1cm}
		\textbf{OPTIONS:}
		\begin{enumerate}
			\item GenAI models are always simple
			\item Opacity, unpredictable output, emergent capabilities, and post-deployment shifts
			\item GenAI models never need governance
			\item GenAI models are always transparent
		\end{enumerate}
		
		\vspace{0.1cm}
		\textcolor{myblue}{\textbf{ANSWER: B}}
		
		\vspace{0.1cm}
		\textbf{DETAILED EXPLANATION:}
		\begin{itemize}
			\item \textbf{Why B is correct:} GenAI governance is harder due to opacity (functioning is almost always opaque), unpredictable output (probabilistic nature), emergent capabilities (behaviors not anticipated), and post-deployment shifts (model behavior can change after deployment).
			\item \textbf{Why A is wrong:} GenAI models are complex, not simple
			\item \textbf{Why C is wrong:} GenAI models require significant governance
			\item \textbf{Why D is wrong:} GenAI models are often opaque
		\end{itemize}
		
		\textcolor{myred}{\textbf{EXAM TRAP:}} GenAI governance is harder due to \textcolor{myblue}{opacity, unpredictability, and emergent capabilities}!
	\end{frame}
	
	\begin{frame}{Q77: Emergent Capabilities - Tutorial}
		\fontsize{6.5}{7.5}\selectfont
		\textbf{QUESTION: What are emergent capabilities in GenAI?}
		
		\vspace{0.1cm}
		\textbf{OPTIONS:}
		\begin{enumerate}
			\item Capabilities that are always predictable
			\item Behaviors that were not anticipated by developers as models become larger and more sophisticated
			\item Capabilities that never change
			\item Capabilities that are always fully understood
		\end{enumerate}
		
		\vspace{0.1cm}
		\textcolor{myblue}{\textbf{ANSWER: B}}
		
		\vspace{0.1cm}
		\textbf{DETAILED EXPLANATION:}
		\begin{itemize}
			\item \textbf{Why B is correct:} The unpredictable nature of output from GenAI systems can lead to emergent capabilities, namely behaviors that were not anticipated by developers. As models become larger and more sophisticated, they begin to display skills that were not observed in smaller models.
			\item \textbf{Why A is wrong:} Emergent capabilities are UNPREDICTABLE
			\item \textbf{Why C is wrong:} Capabilities can change
			\item \textbf{Why D is wrong:} Emergent capabilities may not be fully understood
		\end{itemize}
		
		\textcolor{myred}{\textbf{EXAM TRAP:}} Emergent capabilities are \textcolor{myblue}{unexpected behaviors not anticipated by developers}!
	\end{frame}
	
	\begin{frame}{Q78: Internal Governance Levers - Tutorial}
		\fontsize{6.5}{7.5}\selectfont
		\textbf{QUESTION: What is an internal governance lever for GenAI?}
		
		\vspace{0.1cm}
		\textbf{OPTIONS:}
		\begin{enumerate}
			\item Only relying on external regulation
			\item Assigning ownership and accountability, risk-based review processes, incentivizing responsible behavior, and training staff
			\item Ignoring all internal controls
			\item Only using vendor-provided governance
		\end{enumerate}
		
		\vspace{0.1cm}
		\textcolor{myblue}{\textbf{ANSWER: B}}
		
		\vspace{0.1cm}
		\textbf{DETAILED EXPLANATION:}
		\begin{itemize}
			\item \textbf{Why B is correct:} Internal governance levers include: assigning ownership and accountability across the lifecycle, implementing risk-based review and deployment processes, incentivizing responsible behavior through risk-tied performance metrics, training staff to recognize and escalate GenAI-related risks, and establishing clear escalation paths.
			\item \textbf{Why A is wrong:} External regulation is not an internal lever
			\item \textbf{Why C is wrong:} Internal controls are essential
			\item \textbf{Why D is wrong:} Internal governance is needed beyond vendor solutions
		\end{itemize}
		
		\textcolor{myred}{\textbf{EXAM TRAP:}} Internal governance levers include \textcolor{myblue}{ownership, risk-based reviews, incentives, and training}!
	\end{frame}
	
	\begin{frame}{Q79: External Governance Constraints - Tutorial}
		\fontsize{6.5}{7.5}\selectfont
		\textbf{QUESTION: What is an external governance constraint for GenAI?}
		
		\vspace{0.1cm}
		\textbf{OPTIONS:}
		\begin{enumerate}
			\item Only internal policies
			\item Tracking evolving legal and regulatory requirements and aligning with emerging standards
			\item Ignoring all regulations
			\item Only using vendor governance
		\end{enumerate}
		
		\vspace{0.1cm}
		\textcolor{myblue}{\textbf{ANSWER: B}}
		
		\vspace{0.1cm}
		\textbf{DETAILED EXPLANATION:}
		\begin{itemize}
			\item \textbf{Why B is correct:} External governance constraints include: tracking evolving legal and regulatory requirements (EU AI Act, DSA, US state laws), evaluating and negotiating responsibilities with model providers, and aligning with emerging standards and certifications (NIST AI RMF, ISO/IEC 42001).
			\item \textbf{Why A is wrong:} Internal policies are internal levers
			\item \textbf{Why C is wrong:} Regulations cannot be ignored
			\item \textbf{Why D is wrong:} Vendor governance is just one aspect
		\end{itemize}
		
		\textcolor{myred}{\textbf{EXAM TRAP:}} External governance requires \textcolor{myblue}{tracking regulation and aligning with standards}!
	\end{frame}
	
	\begin{frame}{Q80: Adaptive Tools for GenAI - Tutorial}
		\fontsize{6.5}{7.5}\selectfont
		\textbf{QUESTION: What is an adaptive tool for GenAI governance?}
		
		\vspace{0.1cm}
		\textbf{OPTIONS:}
		\begin{enumerate}
			\item Only one-time reviews
			\item Continuously monitoring post-deployment usage, red-teaming, scenario planning, and feedback loops
			\item Ignoring post-deployment behavior
			\item Only initial validation
		\end{enumerate}
		
		\vspace{0.1cm}
		\textcolor{myblue}{\textbf{ANSWER: B}}
		
		\vspace{0.1cm}
		\textbf{DETAILED EXPLANATION:}
		\begin{itemize}
			\item \textbf{Why B is correct:} Adaptive tools and practices include: continuously monitoring post-deployment usage and system behavior, using red-teaming and adversarial testing to uncover edge cases, developing scenario plans for unpredictability and misuse, deploying technical safeguards, and enabling feedback loops and update mechanisms.
			\item \textbf{Why A is wrong:} One-time reviews are insufficient
			\item \textbf{Why C is wrong:} Post-deployment behavior must be monitored
			\item \textbf{Why D is wrong:} Ongoing monitoring is essential
		\end{itemize}
		
		\textcolor{myred}{\textbf{EXAM TRAP:}} Adaptive tools require \textcolor{myblue}{continuous monitoring, testing, and feedback}!
	\end{frame}
	
	\begin{frame}{Q81: Red-Teaming for GenAI - Tutorial}
		\fontsize{6.5}{7.5}\selectfont
		\textbf{QUESTION: What is the purpose of red-teaming for GenAI?}
		
		\vspace{0.1cm}
		\textbf{OPTIONS:}
		\begin{enumerate}
			\item To make the model faster
			\item To uncover edge cases and vulnerabilities through adversarial testing
			\item To increase model accuracy
			\item To reduce model size
		\end{enumerate}
		
		\vspace{0.1cm}
		\textcolor{myblue}{\textbf{ANSWER: B}}
		
		\vspace{0.1cm}
		\textbf{DETAILED EXPLANATION:}
		\begin{itemize}
			\item \textbf{Why B is correct:} Red-teaming and adversarial testing are used to uncover edge cases, vulnerabilities, and unexpected behaviors. This involves attempting to cause the system to fail or behave in unintended ways.
			\item \textbf{Why A is wrong:} Speed is not the purpose
			\item \textbf{Why C is wrong:} Accuracy is not the primary purpose
			\item \textbf{Why D is wrong:} Model size is not the focus
		\end{itemize}
		
		\textcolor{myred}{\textbf{EXAM TRAP:}} Red-teaming uncovers \textcolor{myblue}{edge cases and vulnerabilities} through adversarial testing!
	\end{frame}
	
	\begin{frame}{Q82: Scenario Planning for GenAI - Tutorial}
		\fontsize{6.5}{7.5}\selectfont
		\textbf{QUESTION: Why is scenario planning important for GenAI governance?}
		
		\vspace{0.1cm}
		\textbf{OPTIONS:}
		\begin{enumerate}
			\item Because GenAI is always predictable
			\item To prepare for unpredictability and misuse by understanding what unpredictable behaviors are most likely and impactful
			\item Because scenario planning is never needed
			\item Because GenAI never fails
		\end{enumerate}
		
		\vspace{0.1cm}
		\textcolor{myblue}{\textbf{ANSWER: B}}
		
		\vspace{0.1cm}
		\textbf{DETAILED EXPLANATION:}
		\begin{itemize}
			\item \textbf{Why B is correct:} Understand what unpredictable behavior and misuse cases are most likely/impactful for your particular use case and develop appropriate scenario plans. Create different incident scenarios, assign roles, and run tabletop exercises to test readiness.
			\item \textbf{Why A is wrong:} GenAI is often unpredictable
			\item \textbf{Why C is wrong:} Scenario planning is essential
			\item \textbf{Why D is wrong:} GenAI can fail in unexpected ways
		\end{itemize}
		
		\textcolor{myred}{\textbf{EXAM TRAP:}} Scenario planning prepares for \textcolor{myblue}{unpredictability and misuse}!
	\end{frame}
	
	\begin{frame}{Q83: Questions for Risk Professionals - Tutorial}
		\fontsize{6.5}{7.5}\selectfont
		\textbf{QUESTION: What is a key question for strategic risk assessment of GenAI?}
		
		\vspace{0.1cm}
		\textbf{OPTIONS:}
		\begin{enumerate}
			\item "What is the fastest way to deploy?"
			\item "Is GenAI the right tool for this application?"
			\item "How can we minimize costs?"
			\item "How can we ignore all risks?"
		\end{enumerate}
		
		\vspace{0.1cm}
		\textcolor{myblue}{\textbf{ANSWER: B}}
		
		\vspace{0.1cm}
		\textbf{DETAILED EXPLANATION:}
		\begin{itemize}
			\item \textbf{Why B is correct:} Key strategic risk assessment questions include: "Is GenAI the right tool for this application? Ensure the reasons for deploying GenAI tools are good reasons and know what trade-offs are to be expected."
			\item \textbf{Why A is wrong:} Speed should not override strategic assessment
			\item \textbf{Why C is wrong:} Cost is a consideration but not the primary question
			\item \textbf{Why D is wrong:} Risks must be addressed, not ignored
		\end{itemize}
		
		\textcolor{myred}{\textbf{EXAM TRAP:}} Strategic assessment asks \textcolor{myblue}{whether GenAI is the right tool} - not just how to deploy quickly!
	\end{frame}
	
	\begin{frame}{Q84: Technical Risk Control Questions - Tutorial}
		\fontsize{6.5}{7.5}\selectfont
		\textbf{QUESTION: What is a key technical risk control question for GenAI?}
		
		\vspace{0.1cm}
		\textbf{OPTIONS:}
		\begin{enumerate}
			\item "Have we tested for unpredictability?"
			\item "How can we make it faster?"
			\item "How can we reduce costs?"
			\item "Can we ignore all safeguards?"
		\end{enumerate}
		
		\vspace{0.1cm}
		\textcolor{myblue}{\textbf{ANSWER: A}}
		
		\vspace{0.1cm}
		\textbf{DETAILED EXPLANATION:}
		\begin{itemize}
			\item \textbf{Why A is correct:} Key technical risk control questions include: "Have we tested for unpredictability? (e.g., hallucinations, model drift, etc.)" and "What safeguards are in place against misuse? (e.g., guardrails, human review triggers, usage limits, logs, etc.)"
			\item \textbf{Why B is wrong:} Speed is not the primary technical control concern
			\item \textbf{Why C is wrong:} Cost is a consideration but not primary
			\item \textbf{Why D is wrong:} Safeguards are essential
		\end{itemize}
		
		\textcolor{myred}{\textbf{EXAM TRAP:}} Technical risk controls focus on \textcolor{myblue}{testing for unpredictability and safeguards}!
	\end{frame}
	
	\begin{frame}{Q85: Organizational Questions for GenAI - Tutorial}
		\fontsize{6.5}{7.5}\selectfont
		\textbf{QUESTION: What is a key organizational question for GenAI governance?}
		
		\vspace{0.1cm}
		\textbf{OPTIONS:}
		\begin{enumerate}
			\item "Who owns risk for GenAI?"
			\item "How can we ignore all risks?"
			\item "Who can benefit the most?"
			\item "How can we avoid responsibility?"
		\end{enumerate}
		
		\vspace{0.1cm}
		\textcolor{myblue}{\textbf{ANSWER: A}}
		
		\vspace{0.1cm}
		\textbf{DETAILED EXPLANATION:}
		\begin{itemize}
			\item \textbf{Why A is correct:} Key organizational questions include: "Who owns risk for GenAI? Is that person/team empowered? Embedding accountability is key." and "Are the staff using GenAI appropriately trained and know of unpredictable behavior?"
			\item \textbf{Why B is wrong:} Risks cannot be ignored
			\item \textbf{Why C is wrong:} Benefit is not the primary governance question
			\item \textbf{Why D is wrong:} Responsibility cannot be avoided
		\end{itemize}
		
		\textcolor{myred}{\textbf{EXAM TRAP:}} Organizational governance requires \textcolor{myblue}{clear ownership and accountability} for GenAI risks!
	\end{frame}
	
	% ============================================================
	% SECTION 8: CHAPTER QUESTIONS - QUESTIONS 86-100
	% ============================================================
	
	\begin{frame}{Q86: Model Development vs Validation - Tutorial}
		\fontsize{6.5}{7.5}\selectfont
		\textbf{QUESTION: When developing and testing a model, should the developer calibrate the model against the full dataset to obtain the best fit?}
		
		\vspace{0.1cm}
		\textbf{OPTIONS:}
		\begin{enumerate}
			\item True - always use the full dataset
			\item False - a holdout sample should be retained for validation
			\item True - calibration should use all available data
			\item False - calibration is not important
		\end{enumerate}
		
		\vspace{0.1cm}
		\textcolor{myblue}{\textbf{ANSWER: B}}
		
		\vspace{0.1cm}
		\textbf{DETAILED EXPLANATION:}
		\begin{itemize}
			\item \textbf{Why B is correct:} Model development should not use the full dataset for calibration. A holdout sample should be retained for validation to assess generalization performance and avoid overfitting.
			\item \textbf{Why A is wrong:} Using all data for calibration prevents proper validation
			\item \textbf{Why C is wrong:} This would prevent proper validation
			\item \textbf{Why D is wrong:} Calibration is important but should use appropriate data splits
		\end{itemize}
		
		\textcolor{myred}{\textbf{EXAM TRAP:}} Always retain \textcolor{myblue}{holdout data for validation} - don't calibrate on full dataset!
	\end{frame}
	
	\begin{frame}{Q87: Model Validation Independence - Tutorial}
		\fontsize{6.5}{7.5}\selectfont
		\textbf{QUESTION: Is model validation part of the model development process?}
		
		\vspace{0.1cm}
		\textbf{OPTIONS:}
		\begin{enumerate}
			\item True - validation is part of development
			\item False - model validation is a separate process that occurs independently of, and generally following, the model development process
			\item True - they are the same thing
			\item False - validation never occurs
		\end{enumerate}
		
		\vspace{0.1cm}
		\textcolor{myblue}{\textbf{ANSWER: B}}
		
		\vspace{0.1cm}
		\textbf{DETAILED EXPLANATION:}
		\begin{itemize}
			\item \textbf{Why B is correct:} Model validation is a separate process that occurs independently of, and generally following, the model development process. It provides independent assessment of model performance.
			\item \textbf{Why A is wrong:} Validation is separate from development
			\item \textbf{Why C is wrong:} They are distinct processes
			\item \textbf{Why D is wrong:} Validation is essential
		\end{itemize}
		
		\textcolor{myred}{\textbf{EXAM TRAP:}} Validation is \textcolor{myblue}{separate from and follows development} - provides independent assessment!
	\end{frame}
	
	\begin{frame}{Q88: Testing Plan Components - Tutorial}
		\fontsize{6.5}{7.5}\selectfont
		\textbf{QUESTION: A good model testing plan includes system integration tests, performance tests, and acceptance tests. Is this true?}
		
		\vspace{0.1cm}
		\textbf{OPTIONS:}
		\begin{enumerate}
			\item True
			\item False
			\item Only system tests are needed
			\item Only acceptance tests are needed
		\end{enumerate}
		
		\vspace{0.1cm}
		\textcolor{myblue}{\textbf{ANSWER: A}}
		
		\vspace{0.1cm}
		\textbf{DETAILED EXPLANATION:}
		\begin{itemize}
			\item \textbf{Why A is correct:} Steps in the testing process include unit tests, component tests, integration tests, system tests as well as system integration tests, performance tests, regression tests, and user acceptance tests. A good testing plan includes multiple types of tests.
			\item \textbf{Why B is wrong:} All these tests are important components of a testing plan
			\item \textbf{Why C is wrong:} Multiple test types are needed
			\item \textbf{Why D is wrong:} Multiple test types are needed
		\end{itemize}
		
		\textcolor{myred}{\textbf{EXAM TRAP:}} A comprehensive testing plan includes \textcolor{myblue}{multiple test types} - not just one!
	\end{frame}
	
	\begin{frame}{Q89: Testing Extreme Cases - Tutorial}
		\fontsize{6.5}{7.5}\selectfont
		\textbf{QUESTION: If a model is considered a 'black box', should unrealistic and extreme test cases be tested?}
		
		\vspace{0.1cm}
		\textbf{OPTIONS:}
		\begin{enumerate}
			\item True - extreme cases should be tested for all models
			\item False - black box models only need normal case testing
			\item True - but only for white box models
			\item False - extreme cases are never useful
		\end{enumerate}
		
		\vspace{0.1cm}
		\textcolor{myblue}{\textbf{ANSWER: A}}
		
		\vspace{0.1cm}
		\textbf{DETAILED EXPLANATION:}
		\begin{itemize}
			\item \textbf{Why A is correct:} Whether a model is white box, gray box or black box, the less realistic the parameters and input data (e.g., missing trades, excessive trades, incorrect scales or mappings, extreme parameter values, etc.), the better. Extreme cases help identify vulnerabilities.
			\item \textbf{Why B is wrong:} Black box models also need extreme case testing
			\item \textbf{Why C is wrong:} All models benefit from extreme case testing
			\item \textbf{Why D is wrong:} Extreme cases are very useful
		\end{itemize}
		
		\textcolor{myred}{\textbf{EXAM TRAP:}} \textcolor{myblue}{Extreme cases should be tested for ALL models} - black box, white box, and gray box!
	\end{frame}
	
	\begin{frame}{Q90: Model Validation Responsibilities - Tutorial}
		\fontsize{6.5}{7.5}\selectfont
		\textbf{QUESTION: Is the model validation team responsible for data selection and testing parameter stability?}
		
		\vspace{0.1cm}
		\textbf{OPTIONS:}
		\begin{enumerate}
			\item True - validation team does everything
			\item False - developers are responsible for data selection and parameter stability testing; validation ensures these have been done appropriately
			\item True - validation team selects all data
			\item False - no one is responsible
		\end{enumerate}
		
		\vspace{0.1cm}
		\textcolor{myblue}{\textbf{ANSWER: B}}
		
		\vspace{0.1cm}
		\textbf{DETAILED EXPLANATION:}
		\begin{itemize}
			\item \textbf{Why B is correct:} Testing the stability of parameters and robustness to various scenarios is performed by the developers, who are also responsible for data selection and specification of any transformations. Model validation just ensures that these have been done appropriately and documented.
			\item \textbf{Why A is wrong:} Developers have primary responsibility for these activities
			\item \textbf{Why C is wrong:} Developers select data
			\item \textbf{Why D is wrong:} Responsibilities are assigned
		\end{itemize}
		
		\textcolor{myred}{\textbf{EXAM TRAP:}} Developers do the work; validation \textcolor{myblue}{ensures it was done appropriately}!
	\end{frame}
	
	\begin{frame}{Q91: Model Results Context - Tutorial}
		\fontsize{6.5}{7.5}\selectfont
		\textbf{QUESTION: Should model results be considered independently, without considering the context of the decision being made?}
		
		\vspace{0.1cm}
		\textbf{OPTIONS:}
		\begin{enumerate}
			\item True - independence ensures objectivity
			\item False - the context of the decision is critical for proper interpretation
			\item True - context is irrelevant
			\item False - results should never be considered
		\end{enumerate}
		
		\vspace{0.1cm}
		\textcolor{myblue}{\textbf{ANSWER: B}}
		
		\vspace{0.1cm}
		\textbf{DETAILED EXPLANATION:}
		\begin{itemize}
			\item \textbf{Why B is correct:} The context of the decision is critical, as misinterpretation can occur when model results are considered independently, without considering the specific context of the decision being made. Risk models should be used as a tool to inform decision making within the appropriate context.
			\item \textbf{Why A is wrong:} Context is essential for proper interpretation
			\item \textbf{Why C is wrong:} Context is highly relevant
			\item \textbf{Why D is wrong:} Results should be considered with context
		\end{itemize}
		
		\textcolor{myred}{\textbf{EXAM TRAP:}} Model results must be interpreted \textcolor{myblue}{within decision context} - not independently!
	\end{frame}
	
	\begin{frame}{Q92: Approximations in VaR - Tutorial}
		\fontsize{6.5}{7.5}\selectfont
		\textbf{QUESTION: A firm's VaR model runs overnight. The trading head says accuracy is more important than speed, so no approximation should be used. Is this correct?}
		
		\vspace{0.1cm}
		\textbf{OPTIONS:}
		\begin{enumerate}
			\item True - accuracy is always more important
			\item False - VaR calculations are not as time sensitive as pricing models, and numerical approximations are often employed
			\item True - approximations should never be used
			\item False - speed is always more important
		\end{enumerate}
		
		\vspace{0.1cm}
		\textcolor{myblue}{\textbf{ANSWER: B}}
		
		\vspace{0.1cm}
		\textbf{DETAILED EXPLANATION:}
		\begin{itemize}
			\item \textbf{Why B is correct:} VaR calculations are not as time sensitive as pricing models used for intraday trading, and numerical approximations are often employed. For example, scaling up the time horizon is frequently used.
			\item \textbf{Why A is wrong:} VaR can use appropriate approximations
			\item \textbf{Why C is wrong:} Approximations can be appropriate for VaR
			\item \textbf{Why D is wrong:} Speed and accuracy must be balanced
		\end{itemize}
		
		\textcolor{myred}{\textbf{EXAM TRAP:}} VaR calculations can use \textcolor{myblue}{appropriate approximations} - not as time sensitive as intraday trading!
	\end{frame}
	
	\begin{frame}{Q93: Model Adaptation vs System Adaptation - Tutorial}
		\fontsize{6.5}{7.5}\selectfont
		\textbf{QUESTION: When model validation detects outdated assumptions and the developer says a rebuild won't work on existing infrastructure, what adaptation is needed?}
		
		\vspace{0.1cm}
		\textbf{OPTIONS:}
		\begin{enumerate}
			\item System adaptation
			\item Core adaptation
			\item Model design
			\item No adaptation is needed
		\end{enumerate}
		
		\vspace{0.1cm}
		\textcolor{myblue}{\textbf{ANSWER: B}}
		
		\vspace{0.1cm}
		\textbf{DETAILED EXPLANATION:}
		\begin{itemize}
			\item \textbf{Why B is correct:} Core adaptation is the process where the model core, and perhaps its interfaces, are adapted (or even replaced). System adaptation describes a situation where the system around the model core is adapted. In this case, a core adaptation is first required.
			\item \textbf{Why A is wrong:} System adaptation is for the system around the core
			\item \textbf{Why C is wrong:} Design is the initial creation phase
			\item \textbf{Why D is wrong:} Adaptation is clearly needed
		\end{itemize}
		
		\textcolor{myred}{\textbf{EXAM TRAP:}} Core adaptation = \textcolor{myblue}{model core changes}; System adaptation = \textcolor{myblue}{system around the core changes}!
	\end{frame}
	
	\begin{frame}{Q94: Core Adaptation Example - Tutorial}
		\fontsize{6.5}{7.5}\selectfont
		\textbf{QUESTION: In the crypto portfolio adaptation example, what is the best implementation task?}
		
		\vspace{0.1cm}
		\textbf{OPTIONS:}
		\begin{enumerate}
			\item Model design
			\item Core adaptation
			\item Model development
			\item System adaptation
		\end{enumerate}
		
		\vspace{0.1cm}
		\textcolor{myblue}{\textbf{ANSWER: B}}
		
		\vspace{0.1cm}
		\textbf{DETAILED EXPLANATION:}
		\begin{itemize}
			\item \textbf{Why B is correct:} Core adaptation is the process where the model core, and perhaps its interfaces, are adapted (or even replaced). Since the existing code can mainly be reused with adaptations and a new data feed added to the interface, core adaptation is the best choice.
			\item \textbf{Why A is wrong:} The model already exists - it needs adaptation
			\item \textbf{Why C is wrong:} Development would be starting from scratch
			\item \textbf{Why D is wrong:} The user interface is not changing
		\end{itemize}
		
		\textcolor{myred}{\textbf{EXAM TRAP:}} Core adaptation = \textcolor{myblue}{reusing existing code with adaptations}!
	\end{frame}
	
	\begin{frame}{Q95: Data Governance Violation - Tutorial}
		\fontsize{6.5}{7.5}\selectfont
		\textbf{QUESTION: A bank sends customer data including SSNs to a vendor for AML checks. Is this a data governance violation?}
		
		\vspace{0.1cm}
		\textbf{OPTIONS:}
		\begin{enumerate}
			\item True - it may be a violation without proper approval
			\item False - it's completely acceptable
			\item True - but only if the data is used incorrectly
			\item False - vendors can always access customer data
		\end{enumerate}
		
		\vspace{0.1cm}
		\textcolor{myblue}{\textbf{ANSWER: A}}
		
		\vspace{0.1cm}
		\textbf{DETAILED EXPLANATION:}
		\begin{itemize}
			\item \textbf{Why A is correct:} The bank needs to be extremely careful when sending non-publicly available customer information to anyone, even their own vendors. Approval from the bank's data governance committee must be sought, and specification of whether the data is for single or recurring use must also be included.
			\item \textbf{Why B is wrong:} Approval is required before sending sensitive data
			\item \textbf{Why C is wrong:} Sending data without proper approval is a violation regardless of use
			\item \textbf{Why D is wrong:} Vendors cannot access customer data without proper controls
		\end{itemize}
		
		\textcolor{myred}{\textbf{EXAM TRAP:}} Sending sensitive customer data to vendors requires \textcolor{myblue}{proper approval and governance}!
	\end{frame}
	
	\begin{frame}{Q96: Generative AI Data Privacy - Tutorial}
		\fontsize{6.5}{7.5}\selectfont
		\textbf{QUESTION: A developer uploads a CSV with SSNs and credit data to an online GenAI platform. The developer states the platform is completely private. Is this correct?}
		
		\vspace{0.1cm}
		\textbf{OPTIONS:}
		\begin{enumerate}
			\item True - online platforms are always private
			\item False - users need to be very careful uploading proprietary data as it may be exposed or used in training
			\item True - as long as the developer deletes it later
			\item False - but only if the data is encrypted
		\end{enumerate}
		
		\vspace{0.1cm}
		\textcolor{myblue}{\textbf{ANSWER: B}}
		
		\vspace{0.1cm}
		\textbf{DETAILED EXPLANATION:}
		\begin{itemize}
			\item \textbf{Why B is correct:} Users of online generative AI systems need to be very careful of uploading any proprietary data at all. Anything uploaded to an AI platform has the potential to be exposed at some point or used in training successive models.
			\item \textbf{Why A is wrong:} No online platform is completely private
			\item \textbf{Why C is wrong:} Deletion may not remove data from training sets
			\item \textbf{Why D is wrong:} Even encrypted data may be vulnerable
		\end{itemize}
		
		\textcolor{myred}{\textbf{EXAM TRAP:}} Never upload \textcolor{myblue}{proprietary or sensitive data} to online GenAI platforms!
	\end{frame}
	
	\begin{frame}{Q97: Final Thoughts - Team Qualification - Tutorial}
		\fontsize{6.5}{7.5}\selectfont
		\textbf{QUESTION: What is essential for validation teams to effectively validate AI/ML models?}
		
		\vspace{0.1cm}
		\textbf{OPTIONS:}
		\begin{enumerate}
			\item Only traditional statistical knowledge
			\item Knowledge of AI/ML techniques, hands-on expertise, and familiarity with AI/ML development platforms
			\item Only programming skills
			\item Only business knowledge
		\end{enumerate}
		
		\vspace{0.1cm}
		\textcolor{myblue}{\textbf{ANSWER: B}}
		
		\vspace{0.1cm}
		\textbf{DETAILED EXPLANATION:}
		\begin{itemize}
			\item \textbf{Why B is correct:} The validation team needs knowledge of and hands-on expertise in AI/ML techniques, strong foundations in product, econometric, statistical, and computer science knowledge, and familiarity with AI/ML development and deployment platforms, cloud environments, low-code/no-code environments, CI/CD and MLOps platforms.
			\item \textbf{Why A is wrong:} Traditional knowledge alone is insufficient
			\item \textbf{Why C is wrong:} Programming alone is insufficient
			\item \textbf{Why D is wrong:} Business knowledge alone is insufficient
		\end{itemize}
		
		\textcolor{myred}{\textbf{EXAM TRAP:}} Validation teams need \textcolor{myblue}{diverse skills across AI/ML, statistics, and platforms}!
	\end{frame}
	
	\begin{frame}{Q98: End-to-End Review - Tutorial}
		\fontsize{6.5}{7.5}\selectfont
		\textbf{QUESTION: Why does the entire model validation framework need review for AI/ML models?}
		
		\vspace{0.1cm}
		\textbf{OPTIONS:}
		\begin{enumerate}
			\item Because AI/ML models are not subject to regulation
			\item Because regulators consider AI/ML models fall within the scope of model risk management and SR 11-7
			\item Because AI/ML models never need review
			\item Because traditional frameworks are always sufficient
		\end{enumerate}
		
		\vspace{0.1cm}
		\textcolor{myblue}{\textbf{ANSWER: B}}
		
		\vspace{0.1cm}
		\textbf{DETAILED EXPLANATION:}
		\begin{itemize}
			\item \textbf{Why B is correct:} The entire model validation framework needs to be reviewed as regulators consider AI/ML models fall within the scope of model risk management and SR 11-7, being principle based, is also applicable to AI/ML models.
			\item \textbf{Why A is wrong:} AI/ML models ARE subject to regulation
			\item \textbf{Why C is wrong:} AI/ML models require review
			\item \textbf{Why D is wrong:} Traditional frameworks may need adaptation
		\end{itemize}
		
		\textcolor{myred}{\textbf{EXAM TRAP:}} SR 11-7 applies to \textcolor{myblue}{AI/ML models as well} - need end-to-end review!
	\end{frame}
	
	\begin{frame}{Q99: Model Complexity Balance - Tutorial}
		\fontsize{6.5}{7.5}\selectfont
		\textbf{QUESTION: Is more complex always better for AI/ML models?}
		
		\vspace{0.1cm}
		\textbf{OPTIONS:}
		\begin{enumerate}
			\item True - more complexity always means better performance
			\item False - an appropriate balance needs to be found between model performance and interpretability, cost, and other factors
			\item True - complexity is always desirable
			\item False - simpler models are always better
		\end{enumerate}
		
		\vspace{0.1cm}
		\textcolor{myblue}{\textbf{ANSWER: B}}
		
		\vspace{0.1cm}
		\textbf{DETAILED EXPLANATION:}
		\begin{itemize}
			\item \textbf{Why B is correct:} An appropriate balance needs to be found between model performance and all the other factors (e.g., interpretability, feedback from the supervisor, cost and effort of model implementation, maintenance and monitoring, in-house expertise, etc.).
			\item \textbf{Why A is wrong:} More complexity is not always better
			\item \textbf{Why C is wrong:} Complexity has costs and trade-offs
			\item \textbf{Why D is wrong:} Simpler models are not always better either
		\end{itemize}
		
		\textcolor{myred}{\textbf{EXAM TRAP:}} Balance is needed between \textcolor{myblue}{performance and interpretability, cost, and other factors}!
	\end{frame}
	
	\begin{frame}{Q100: Final Summary - Data and AI Model Governance - Tutorial}
		\fontsize{6.5}{7.5}\selectfont
		\textbf{QUESTION: What is the most comprehensive summary of data and AI model governance?}
		
		\vspace{0.1cm}
		\textbf{OPTIONS:}
		\begin{enumerate}
			\item Only model validation is needed
			\item Data governance, model governance, model development, model validation, three lines of defense, and GenAI governance are all essential components
			\item Only data governance is needed
			\item Only GenAI governance is needed
		\end{enumerate}
		
		\vspace{0.1cm}
		\textcolor{myblue}{\textbf{ANSWER: B}}
		
		\vspace{0.1cm}
		\textbf{DETAILED EXPLANATION:}
		\begin{itemize}
			\item \textbf{Why B is correct:} Comprehensive data and AI model governance requires: Data Governance (strategy, quality, provenance), Model Governance (policies, documentation, inventory), Model Development (objectives, data, training, testing), Model Validation (initial, ongoing, change-based), Three Lines of Defense (owners, risk management, audit), and GenAI Governance (internal levers, external constraints, adaptive tools).
			\item \textbf{Why A is wrong:} Validation is just one component
			\item \textbf{Why C is wrong:} Model governance is also needed
			\item \textbf{Why D is wrong:} All aspects of governance are needed
		\end{itemize}
		
		\textcolor{myred}{\textbf{EXAM SUCCESS:}} Understand \textcolor{myblue}{data governance, model risk management, validation, and GenAI governance}!
	\end{frame}
	
\end{document}